\import{algebra/ring/regular-module.tex}
\import{algebra/ring/homomorphism.tex}
\import{algebra/ring/divisibility.tex}
\import{algebra/ring/power.tex}
\import{algebra/module/submodule.tex}
\import{set/closure-system.tex}

\section{Ideals}

Left and right ideals are submodules of the corresponding regular modules.
Unqualified ideals are two-sided ideals; in a commutative ring this agrees
with either one-sided notion.  The expanded characterizations below expose
the familiar additive-subgroup and absorption conditions for ordinary use.

\begin{definition}\label{left_ideal}
    $I$ is a left ideal of $R$ iff
        $R$ is a ring and
        $I$ is a submodule of $\LeftRegularModule{R}$ over $R$.
\end{definition}

\begin{proposition}\label{left_ideal_iff}
    $I$ is a left ideal of $R$ iff
        $R$ is a ring and
        $I$ is a submodule of $\LeftRegularModule{R}$ over $R$.
\end{proposition}
\begin{proof}
    Follows by \obv{left_ideal}.
\end{proof}

\begin{definition}\label{right_ideal}
    $I$ is a right ideal of $R$ iff
        $R$ is a ring and
        $I$ is a submodule of $\RightRegularModule{R}$ over
            $\OppositeRing{R}$.
\end{definition}

\begin{proposition}\label{right_ideal_iff}
    $I$ is a right ideal of $R$ iff
        $R$ is a ring and
        $I$ is a submodule of $\RightRegularModule{R}$ over
            $\OppositeRing{R}$.
\end{proposition}
\begin{proof}
    Follows by \obv{right_ideal}.
\end{proof}

\begin{definition}\label{two_sided_ideal}
    $I$ is a two-sided ideal of $R$ iff
        $R$ is a ring and
        $I$ is a left ideal of $R$ and
        $I$ is a right ideal of $R$.
\end{definition}

\begin{proposition}\label{two_sided_ideal_iff}
    $I$ is a two-sided ideal of $R$ iff
        $R$ is a ring and
        $I$ is a left ideal of $R$ and
        $I$ is a right ideal of $R$.
\end{proposition}
\begin{proof}
    Follows by \obv{two_sided_ideal}.
\end{proof}

\begin{definition}\label{ideal}
    $I$ is an ideal of $R$ iff
        $I$ is a two-sided ideal of $R$.
\end{definition}

\begin{proposition}\label{ideal_iff}
    $I$ is an ideal of $R$ iff
        $I$ is a two-sided ideal of $R$.
\end{proposition}
\begin{proof}
    Follows by \obv{ideal}.
\end{proof}

\subsection{Guards and expanded characterizations}

\begin{proposition}\label{left_ideal_ring}
    Suppose $I$ is a left ideal of $R$.
    Then $R$ is a ring.
\end{proposition}
\begin{proof}
    Follows by \obv{left_ideal}.
\end{proof}

\begin{proposition}\label{right_ideal_ring}
    Suppose $I$ is a right ideal of $R$.
    Then $R$ is a ring.
\end{proposition}
\begin{proof}
    Follows by \obv{right_ideal}.
\end{proof}

\begin{proposition}\label{two_sided_ideal_ring}
    Suppose $I$ is a two-sided ideal of $R$.
    Then $R$ is a ring.
\end{proposition}
\begin{proof}
    Follows by \obv{two_sided_ideal}.
\end{proof}

\begin{proposition}\label{ideal_ring}
    Suppose $I$ is an ideal of $R$.
    Then $R$ is a ring.
\end{proposition}
\begin{proof}
    $I$ is a two-sided ideal of $R$ by \obv{ideal}.
    Follows by \obv{two_sided_ideal_ring}.
\end{proof}

\begin{proposition}\label{left_ideal_is_submodule}
    Suppose $I$ is a left ideal of $R$.
    Then $I$ is a submodule of $\LeftRegularModule{R}$ over $R$.
\end{proposition}
\begin{proof}
    Follows by \obv{left_ideal}.
\end{proof}

\begin{proposition}\label{right_ideal_is_submodule}
    Suppose $I$ is a right ideal of $R$.
    Then $I$ is a submodule of $\RightRegularModule{R}$ over
        $\OppositeRing{R}$.
\end{proposition}
\begin{proof}
    Follows by \obv{right_ideal}.
\end{proof}

\begin{theorem}\label{left_ideal_expanded_iff}
    $I$ is a left ideal of $R$ iff
        $R$ is a ring and
        $\subgroup{I}{\AdditiveGroup{R}}$ and
        for all $a,x$ such that $a\in\carrier[R]$ and $x\in I$ we have
            $\ringMul[R](a,x)\in I$.
\end{theorem}
\begin{proof}
    Show that if $I$ is a left ideal of $R$, then
        $R$ is a ring and
        $\subgroup{I}{\AdditiveGroup{R}}$ and
        for all $a,x$ such that $a\in\carrier[R]$ and $x\in I$ we have
            $\ringMul[R](a,x)\in I$.
    \begin{subproof}
        Assume $I$ is a left ideal of $R$.
        Hence $R$ is a ring and
            $I$ is a submodule of $\LeftRegularModule{R}$ over $R$ by
            \obv{left_ideal}.
        Hence $\subgroup{I}{
            \ModuleAdditiveGroup{\LeftRegularModule{R}}}$ by
            \obv{submodule_additive_subgroup}.
        Hence $\subgroup{I}{\AdditiveGroup{R}}$ by
            \obv{left_regular_module_additive_group}.
        For all $a,x$ such that $a\in\carrier[R]$ and $x\in I$ we have
            $\scalarMul[\LeftRegularModule{R}](a,x)\in I$ by
            \obv{submodule_scalar}.
        For all $a,x$ such that $a\in\carrier[R]$ and $x\in I$ we have
            $x\in\carrier[R]$ by
            \auto{subgroup_subseteq,additive_group_fields,subseteq}.
        Hence for all $a,x$ such that $a\in\carrier[R]$ and $x\in I$ we have
            $\ringMul[R](a,x)\in I$ by
            \auto{left_regular_scalar_apply}.
    \end{subproof}
    Show that if $R$ is a ring and
        $\subgroup{I}{\AdditiveGroup{R}}$ and
        for all $a,x$ such that $a\in\carrier[R]$ and $x\in I$ we have
            $\ringMul[R](a,x)\in I$, then
        $I$ is a left ideal of $R$.
    \begin{subproof}
        Assume $R$ is a ring and
            $\subgroup{I}{\AdditiveGroup{R}}$ and
            for all $a,x$ such that $a\in\carrier[R]$ and $x\in I$ we have
                $\ringMul[R](a,x)\in I$.
        $\LeftRegularModule{R}$ is a left module over $R$ by
            \obv{left_regular_module_is_exact}.
        $\ModuleAdditiveGroup{\LeftRegularModule{R}}
            =\AdditiveGroup{R}$ by
            \obv{left_regular_module_additive_group}.
        For all $a,x$ such that $a\in\carrier[R]$ and $x\in I$ we have
            $x\in\carrier[R]$ by
            \auto{subgroup_subseteq,additive_group_fields,subseteq}.
        Hence for all $a,x$ such that $a\in\carrier[R]$ and $x\in I$ we have
            $\scalarMul[\LeftRegularModule{R}](a,x)\in I$ by
            \auto{left_regular_scalar_apply}.
        Hence $I$ is a submodule of $\LeftRegularModule{R}$ over $R$ by
            \obv{submodule}.
        Follows by \obv{left_ideal}.
    \end{subproof}
    Follows by \auto.
\end{proof}

\begin{theorem}\label{right_ideal_expanded_iff}
    $I$ is a right ideal of $R$ iff
        $R$ is a ring and
        $\subgroup{I}{\AdditiveGroup{R}}$ and
        for all $x,a$ such that $x\in I$ and $a\in\carrier[R]$ we have
            $\ringMul[R](x,a)\in I$.
\end{theorem}
\begin{proof}
    Show that if $I$ is a right ideal of $R$, then
        $R$ is a ring and
        $\subgroup{I}{\AdditiveGroup{R}}$ and
        for all $x,a$ such that $x\in I$ and $a\in\carrier[R]$ we have
            $\ringMul[R](x,a)\in I$.
    \begin{subproof}
        Assume $I$ is a right ideal of $R$.
        Hence $R$ is a ring and
            $I$ is a submodule of $\RightRegularModule{R}$ over
                $\OppositeRing{R}$ by
            \obv{right_ideal}.
        Hence $\subgroup{I}{
            \ModuleAdditiveGroup{\RightRegularModule{R}}}$ by
            \obv{submodule_additive_subgroup}.
        Hence $\subgroup{I}{\AdditiveGroup{R}}$ by
            \obv{right_regular_module_additive_group}.
        For all $x,a$ such that $x\in I$ and $a\in\carrier[R]$ we have
            $\rightScalarMul{\RightRegularModule{R}}{x}{a}\in I$ by
            \auto{right_scalar_multiplication,submodule_scalar,
                opposite_ring_fields}.
        Hence for all $x,a$ such that $x\in I$ and $a\in\carrier[R]$ we have
            $\ringMul[R](x,a)\in I$ by
            \auto{subgroup_subseteq,additive_group_fields,subseteq,
                right_regular_scalar_apply}.
    \end{subproof}
    Show that if $R$ is a ring and
        $\subgroup{I}{\AdditiveGroup{R}}$ and
        for all $x,a$ such that $x\in I$ and $a\in\carrier[R]$ we have
            $\ringMul[R](x,a)\in I$, then
        $I$ is a right ideal of $R$.
    \begin{subproof}
        Assume $R$ is a ring and
            $\subgroup{I}{\AdditiveGroup{R}}$ and
            for all $x,a$ such that $x\in I$ and $a\in\carrier[R]$ we have
                $\ringMul[R](x,a)\in I$.
        $\RightRegularModule{R}$ is a left module over
            $\OppositeRing{R}$ by
            \obv{right_regular_module_is_exact}.
        $\ModuleAdditiveGroup{\RightRegularModule{R}}
            =\AdditiveGroup{R}$ by
            \obv{right_regular_module_additive_group}.
        For all $a,x$ such that $a\in\carrier[\OppositeRing{R}]$ and
            $x\in I$ we have
            $a\in\carrier[R]$ and $x\in\carrier[R]$ by
            \auto{opposite_ring_fields,subgroup_subseteq,
                additive_group_fields,subseteq}.
        Hence for all $a,x$ such that
            $a\in\carrier[\OppositeRing{R}]$ and $x\in I$ we have
            $\scalarMul[\RightRegularModule{R}](a,x)\in I$ by
            \auto{right_scalar_multiplication,
                right_regular_scalar_apply}.
        Hence $I$ is a submodule of $\RightRegularModule{R}$ over
            $\OppositeRing{R}$ by \obv{submodule}.
        Follows by \obv{right_ideal}.
    \end{subproof}
    Follows by \auto.
\end{proof}

\begin{corollary}\label{left_ideal_additive_subgroup}
    Suppose $I$ is a left ideal of $R$.
    Then $\subgroup{I}{\AdditiveGroup{R}}$.
\end{corollary}
\begin{proof}
    Follows by \auto{left_ideal_expanded_iff}.
\end{proof}

\begin{corollary}\label{right_ideal_additive_subgroup}
    Suppose $I$ is a right ideal of $R$.
    Then $\subgroup{I}{\AdditiveGroup{R}}$.
\end{corollary}
\begin{proof}
    Follows by \auto{right_ideal_expanded_iff}.
\end{proof}

\begin{corollary}\label{left_ideal_subseteq_carrier}
    Suppose $I$ is a left ideal of $R$.
    Then $I\subseteq\carrier[R]$.
\end{corollary}
\begin{proof}
    Follows by
        \auto{left_ideal_additive_subgroup,subgroup_subseteq,
            additive_group_fields}.
\end{proof}

\begin{corollary}\label{right_ideal_subseteq_carrier}
    Suppose $I$ is a right ideal of $R$.
    Then $I\subseteq\carrier[R]$.
\end{corollary}
\begin{proof}
    Follows by
        \auto{right_ideal_additive_subgroup,subgroup_subseteq,
            additive_group_fields}.
\end{proof}

\begin{corollary}\label{two_sided_ideal_subseteq_carrier}
    Suppose $I$ is a two-sided ideal of $R$.
    Then $I\subseteq\carrier[R]$.
\end{corollary}
\begin{proof}
    $I$ is a left ideal of $R$ by \obv{two_sided_ideal}.
    Follows by \obv{left_ideal_subseteq_carrier}.
\end{proof}

\begin{theorem}\label{ideal_finite_sum}
    Suppose $I$ is a two-sided ideal of $R$.
    Suppose $A$ is finite and $f\in\funs{A}{I}$.
    Then $\finiteSum{\AdditiveGroup{R}}{A}{f}\in I$.
\end{theorem}
\begin{proof}
    Let $G=\AdditiveGroup{R}$ and
        $H=\subgroupStructure{G}{I}$.
    Let $i=\subgroupInclusion{G}{I}$.
    $G$ is a commutative group by
        \obv{two_sided_ideal_ring,
            additive_group_is_exact_commutative_group}.
    $\subgroup{I}{G}$ by
        \obv{two_sided_ideal,left_ideal_additive_subgroup}.
    Hence $H$ is a group by \obv{subgroup_structure_is_group}.
    Hence $G$ is a monoid by
        \obv{commutative_group_is_group,group_is_monoid}.
    $H$ is a monoid by \obv{group_is_monoid}.
    $\carrier[H]=I$ by \auto{subgroup_structure_carrier}.
    Hence $f\in\funs{A}{\carrier[H]}$.
    For all $x,y\in\carrier[G]$ we have
        $\mul[G](x,y)=\mul[G](y,x)$ by
        \obv{commutative_group_commutative}.
    Show that for all $x,y\in\carrier[H]$ we have
        $\mul[H](x,y)=\mul[H](y,x)$.
    \begin{subproof}
        Fix $x,y\in\carrier[H]$.
        Hence $x,y\in I$.
        Hence $x,y\in\carrier[G]$ by
            \auto{subgroup_subseteq,subseteq}.
        $\mul[H](x,y)=\mul[G](x,y)$ and
            $\mul[H](y,x)=\mul[G](y,x)$ by
            \auto{subgroup_structure_mul_apply}.
        Follows by \obv{commutative_group_commutative}.
    \end{subproof}
    $\groupHom{i}{H}{G}$ by
        \auto{subgroup_inclusion_group_hom}.
    Hence $i$ is a monoid homomorphism from $H$ to $G$ by
        \obv{group_hom_is_monoid_hom}.
    Hence
        $i(\finiteFold{H}{A}{f})
            =\finiteFold{G}{A}{i\circ f}$ by
        \obv{monoid_hom_preserves_finite_fold}.
    Show that $i\circ f=f$.
    \begin{subproof}
        $i\circ f\in\funs{A}{\carrier[G]}$ by
            \auto{group_hom_type,funs_circ}.
        $I\subseteq\carrier[G]$ by \obv{subgroup_subseteq}.
        Hence $f\in\funs{A}{\carrier[G]}$ by
            \obv{funs_weaken_codom}.
        Show that for every $x\in A$ we have $(i\circ f)(x)=f(x)$.
        \begin{subproof}
            Fix $x\in A$.
            $f(x)\in I$ by \obv{funs_type_apply}.
            $(i\circ f)(x)=i(f(x))$ by
                \auto{circ_apply,group_hom_type,funs_is_function,
                    funs_ran,funs,funs_elim,funs_circ}.
            $i(f(x))=f(x)$ by
                \auto{subgroup_inclusion_apply}.
        \end{subproof}
        Follows by \auto{funs_elim,funext_on_domain}.
    \end{subproof}
    $\finiteFold{H}{A}{f}\in\carrier[H]$ by
        \auto{monoid_finite_fold_canonical}.
    Hence $\finiteFold{H}{A}{f}\in I$.
    $i(\finiteFold{H}{A}{f})=\finiteFold{H}{A}{f}$ by
        \auto{subgroup_inclusion_apply}.
    Hence $\finiteFold{G}{A}{f}\in I$.
    Follows by \auto{finite_sum}.
\end{proof}

\begin{theorem}\label{commutative_left_ideal_iff_right_ideal}
    Let $R$ be a commutative ring.
    Then $I$ is a left ideal of $R$ iff $I$ is a right ideal of $R$.
\end{theorem}
\begin{proof}
    $R$ is a ring by \obv{commutative_ring_is_ring}.
    Show that if $I$ is a left ideal of $R$, then
        $I$ is a right ideal of $R$.
    \begin{subproof}
        Assume $I$ is a left ideal of $R$.
        Hence $\subgroup{I}{\AdditiveGroup{R}}$ by
            \auto{left_ideal_expanded_iff}.
        Show that for all $x,a$ such that $x\in I$ and
            $a\in\carrier[R]$ we have $\ringMul[R](x,a)\in I$.
        \begin{subproof}
            Fix $x,a$.
            Assume $x\in I$ and $a\in\carrier[R]$.
            Hence $x\in\carrier[R]$ by
                \auto{left_ideal_subseteq_carrier,subseteq}.
            $\ringMul[R](x,a)=\ringMul[R](a,x)$ by
                \obv{commutative_ring_mul_commutative}.
            $\ringMul[R](a,x)\in I$ by
                \auto{left_ideal_expanded_iff}.
        \end{subproof}
        Follows by \auto{right_ideal_expanded_iff}.
    \end{subproof}
    Show that if $I$ is a right ideal of $R$, then
        $I$ is a left ideal of $R$.
    \begin{subproof}
        Assume $I$ is a right ideal of $R$.
        Hence $\subgroup{I}{\AdditiveGroup{R}}$ by
            \auto{right_ideal_expanded_iff}.
        Show that for all $a,x$ such that $a\in\carrier[R]$ and
            $x\in I$ we have $\ringMul[R](a,x)\in I$.
        \begin{subproof}
            Fix $a,x$.
            Assume $a\in\carrier[R]$ and $x\in I$.
            Hence $x\in\carrier[R]$ by
                \auto{right_ideal_subseteq_carrier,subseteq}.
            $\ringMul[R](a,x)=\ringMul[R](x,a)$ by
                \obv{commutative_ring_mul_commutative}.
            $\ringMul[R](x,a)\in I$ by
                \auto{right_ideal_expanded_iff}.
        \end{subproof}
        Follows by \auto{left_ideal_expanded_iff}.
    \end{subproof}
    Follows by \auto.
\end{proof}

\begin{theorem}\label{commutative_left_ideal_iff_two_sided_ideal}
    Let $R$ be a commutative ring.
    Then $I$ is a left ideal of $R$ iff
        $I$ is a two-sided ideal of $R$.
\end{theorem}
\begin{proof}
    $R$ is a ring by \obv{commutative_ring_is_ring}.
    $I$ is a left ideal of $R$ iff $I$ is a right ideal of $R$ by
        \obv{commutative_left_ideal_iff_right_ideal}.
    Follows by \auto{two_sided_ideal}.
\end{proof}

\begin{theorem}\label{commutative_ideal_iff_left_ideal}
    Let $R$ be a commutative ring.
    Then $I$ is an ideal of $R$ iff $I$ is a left ideal of $R$.
\end{theorem}
\begin{proof}
    Follows by
        \obv{ideal,commutative_left_ideal_iff_two_sided_ideal}.
\end{proof}

\begin{corollary}\label{commutative_ideal_iff_right_ideal}
    Let $R$ be a commutative ring.
    Then $I$ is an ideal of $R$ iff $I$ is a right ideal of $R$.
\end{corollary}
\begin{proof}
    Follows by
        \auto{commutative_ideal_iff_left_ideal,
            commutative_left_ideal_iff_right_ideal}.
\end{proof}

\begin{corollary}\label{commutative_ideal_iff_two_sided_ideal}
    Let $R$ be a commutative ring.
    Then $I$ is an ideal of $R$ iff $I$ is a two-sided ideal of $R$.
\end{corollary}
\begin{proof}
    Follows by \obv{ideal}.
\end{proof}

\subsection{Zero, unit, and intersections}

\begin{theorem}\label{zero_is_left_ideal}
    Let $R$ be a ring.
    Then $\{\ringZero[R]\}$ is a left ideal of $R$.
\end{theorem}
\begin{proof}
    $\AdditiveGroup{R}$ is a group by
        \obv{additive_group_is_exact_commutative_group,
            commutative_group_is_group}.
    Hence $\subgroup{
        \trivialSubgroupCarrier{\AdditiveGroup{R}}}{\AdditiveGroup{R}}$ by
        \obv{trivial_carrier_subgroup}.
    $\trivialSubgroupCarrier{\AdditiveGroup{R}}=\{\ringZero[R]\}$ by
        \auto{trivial_subgroup_carrier,additive_group_fields}.
    Hence $\subgroup{\{\ringZero[R]\}}{\AdditiveGroup{R}}$.
    For all $a,x$ such that $a\in\carrier[R]$ and
        $x\in\{\ringZero[R]\}$ we have
        $\ringMul[R](a,x)=\ringZero[R]$ by
        \auto{singleton_iff,ring_mul_zero_right}.
    Follows by \auto{left_ideal_expanded_iff,singleton_intro}.
\end{proof}

\begin{theorem}\label{zero_is_right_ideal}
    Let $R$ be a ring.
    Then $\{\ringZero[R]\}$ is a right ideal of $R$.
\end{theorem}
\begin{proof}
    $\subgroup{\{\ringZero[R]\}}{\AdditiveGroup{R}}$ by
        \obv{zero_is_left_ideal,left_ideal_additive_subgroup}.
    For all $x,a$ such that $x\in\{\ringZero[R]\}$ and
        $a\in\carrier[R]$ we have
        $\ringMul[R](x,a)=\ringZero[R]$ by
        \auto{singleton_iff,ring_mul_zero_left}.
    Follows by \auto{right_ideal_expanded_iff,singleton_intro}.
\end{proof}

\begin{theorem}\label{zero_is_two_sided_ideal}
    Let $R$ be a ring.
    Then $\{\ringZero[R]\}$ is a two-sided ideal of $R$.
\end{theorem}
\begin{proof}
    Follows by
        \obv{two_sided_ideal,zero_is_left_ideal,zero_is_right_ideal}.
\end{proof}

\begin{corollary}\label{zero_is_ideal}
    Let $R$ be a ring.
    Then $\{\ringZero[R]\}$ is an ideal of $R$.
\end{corollary}
\begin{proof}
    Follows by \obv{ideal,zero_is_two_sided_ideal}.
\end{proof}

\begin{theorem}\label{unit_is_left_ideal}
    Let $R$ be a ring.
    Then $\carrier[R]$ is a left ideal of $R$.
\end{theorem}
\begin{proof}
    $\AdditiveGroup{R}$ is a group by
        \obv{additive_group_is_exact_commutative_group,
            commutative_group_is_group}.
    $\subgroup{\carrier[\AdditiveGroup{R}]}{\AdditiveGroup{R}}$ by
        \obv{full_carrier_subgroup}.
    Hence $\subgroup{\carrier[R]}{\AdditiveGroup{R}}$ by
        \auto{additive_group_fields}.
    Follows by \auto{left_ideal_expanded_iff,ring_mul_type}.
\end{proof}

\begin{theorem}\label{unit_is_right_ideal}
    Let $R$ be a ring.
    Then $\carrier[R]$ is a right ideal of $R$.
\end{theorem}
\begin{proof}
    $\subgroup{\carrier[R]}{\AdditiveGroup{R}}$ by
        \obv{unit_is_left_ideal,left_ideal_additive_subgroup}.
    Follows by \auto{right_ideal_expanded_iff,ring_mul_type}.
\end{proof}

\begin{theorem}\label{unit_is_two_sided_ideal}
    Let $R$ be a ring.
    Then $\carrier[R]$ is a two-sided ideal of $R$.
\end{theorem}
\begin{proof}
    Follows by
        \obv{two_sided_ideal,unit_is_left_ideal,unit_is_right_ideal}.
\end{proof}

\begin{corollary}\label{unit_is_ideal}
    Let $R$ be a ring.
    Then $\carrier[R]$ is an ideal of $R$.
\end{corollary}
\begin{proof}
    Follows by \obv{ideal,unit_is_two_sided_ideal}.
\end{proof}

\begin{theorem}\label{left_ideal_contains_one_iff_full}
    Suppose $I$ is a left ideal of $R$.
    Then $\ringOne[R]\in I$ iff $I=\carrier[R]$.
\end{theorem}
\begin{proof}
    Show that if $\ringOne[R]\in I$, then $I=\carrier[R]$.
    \begin{subproof}
        Assume $\ringOne[R]\in I$.
        Show that $\carrier[R]\subseteq I$.
        \begin{subproof}
            Show that for every $x\in\carrier[R]$ we have $x\in I$.
            \begin{subproof}
                Fix $x\in\carrier[R]$.
                $\ringMul[R](x,\ringOne[R])\in I$ by
                    \auto{left_ideal_expanded_iff}.
                Follows by \obv{ring_mul_one_right}.
            \end{subproof}
            Follows by \obv{subseteq}.
        \end{subproof}
        $I\subseteq\carrier[R]$ by
            \obv{left_ideal_subseteq_carrier}.
        Follows by \obv{subseteq_antisymmetric}.
    \end{subproof}
    Show that if $I=\carrier[R]$, then $\ringOne[R]\in I$.
    \begin{subproof}
        Assume $I=\carrier[R]$.
        Follows by \obv{left_ideal_ring,ring_one_type}.
    \end{subproof}
    Follows by \auto.
\end{proof}

\begin{theorem}\label{right_ideal_contains_one_iff_full}
    Suppose $I$ is a right ideal of $R$.
    Then $\ringOne[R]\in I$ iff $I=\carrier[R]$.
\end{theorem}
\begin{proof}
    Show that if $\ringOne[R]\in I$, then $I=\carrier[R]$.
    \begin{subproof}
        Assume $\ringOne[R]\in I$.
        Show that $\carrier[R]\subseteq I$.
        \begin{subproof}
            Show that for every $x\in\carrier[R]$ we have $x\in I$.
            \begin{subproof}
                Fix $x\in\carrier[R]$.
                $\ringMul[R](\ringOne[R],x)\in I$ by
                    \auto{right_ideal_expanded_iff}.
                Follows by \obv{ring_mul_one_left}.
            \end{subproof}
            Follows by \obv{subseteq}.
        \end{subproof}
        $I\subseteq\carrier[R]$ by
            \obv{right_ideal_subseteq_carrier}.
        Follows by \obv{subseteq_antisymmetric}.
    \end{subproof}
    Show that if $I=\carrier[R]$, then $\ringOne[R]\in I$.
    \begin{subproof}
        Assume $I=\carrier[R]$.
        Follows by \obv{right_ideal_ring,ring_one_type}.
    \end{subproof}
    Follows by \auto.
\end{proof}

\begin{corollary}\label{left_ideal_proper_iff}
    Suppose $I$ is a left ideal of $R$.
    Then $I\neq\carrier[R]$ iff $\ringOne[R]\notin I$.
\end{corollary}
\begin{proof}
    Follows by \auto{left_ideal_contains_one_iff_full}.
\end{proof}

\begin{corollary}\label{right_ideal_proper_iff}
    Suppose $I$ is a right ideal of $R$.
    Then $I\neq\carrier[R]$ iff $\ringOne[R]\notin I$.
\end{corollary}
\begin{proof}
    Follows by \auto{right_ideal_contains_one_iff_full}.
\end{proof}

\begin{theorem}\label{left_ideal_intersection}
    Suppose $I$ is a left ideal of $R$ and $J$ is a left ideal of $R$.
    Then $I\inter J$ is a left ideal of $R$.
\end{theorem}
\begin{proof}
    $I$ and $J$ are submodules of $\LeftRegularModule{R}$ over $R$ by
        \obv{left_ideal_is_submodule}.
    Hence $I\inter J$ is a submodule of $\LeftRegularModule{R}$ over $R$
        by \obv{submodule_intersection}.
    Follows by \obv{left_ideal}.
\end{proof}

\begin{theorem}\label{right_ideal_intersection}
    Suppose $I$ is a right ideal of $R$ and $J$ is a right ideal of $R$.
    Then $I\inter J$ is a right ideal of $R$.
\end{theorem}
\begin{proof}
    $I$ and $J$ are submodules of $\RightRegularModule{R}$ over
        $\OppositeRing{R}$ by
        \obv{right_ideal_is_submodule}.
    Hence $I\inter J$ is a submodule of $\RightRegularModule{R}$ over
        $\OppositeRing{R}$ by \obv{submodule_intersection}.
    Follows by \obv{right_ideal}.
\end{proof}

\begin{theorem}\label{two_sided_ideal_intersection}
    Suppose $I$ is a two-sided ideal of $R$ and
        $J$ is a two-sided ideal of $R$.
    Then $I\inter J$ is a two-sided ideal of $R$.
\end{theorem}
\begin{proof}
    Follows by
        \obv{two_sided_ideal,left_ideal_intersection,
            right_ideal_intersection}.
\end{proof}

\begin{corollary}\label{ideal_intersection}
    Suppose $I$ is an ideal of $R$ and $J$ is an ideal of $R$.
    Then $I\inter J$ is an ideal of $R$.
\end{corollary}
\begin{proof}
    $I$ and $J$ are two-sided ideals of $R$ by \obv{ideal}.
    Hence $I\inter J$ is a two-sided ideal of $R$ by
        \obv{two_sided_ideal_intersection}.
    Follows by \obv{ideal}.
\end{proof}

\begin{theorem}\label{left_ideal_indexed_intersection}
    Let $R$ be a ring.
    Let $F\in\funs{K}{\pow{\carrier[R]}}$.
    Suppose for every $k\in K$ we have $F(k)$ is a left ideal of $R$.
    Then $\indexedIntersection{K}{F}{\carrier[R]}$
        is a left ideal of $R$.
\end{theorem}
\begin{proof}
    $\LeftRegularModule{R}$ is a left module over $R$ by
        \obv{left_regular_module_is_exact}.
    $\carrier[\LeftRegularModule{R}]=\carrier[R]$ by
        \obv{left_regular_module_carrier}.
    For every $k\in K$ we have
        $F(k)$ is a submodule of $\LeftRegularModule{R}$ over $R$ by
        \obv{left_ideal_is_submodule}.
    Hence $\indexedIntersection{K}{F}{\carrier[R]}$
        is a submodule of $\LeftRegularModule{R}$ over $R$ by
        \auto{submodule_indexed_intersection}.
    Follows by \obv{left_ideal}.
\end{proof}

\begin{theorem}\label{right_ideal_indexed_intersection}
    Let $R$ be a ring.
    Let $F\in\funs{K}{\pow{\carrier[R]}}$.
    Suppose for every $k\in K$ we have $F(k)$ is a right ideal of $R$.
    Then $\indexedIntersection{K}{F}{\carrier[R]}$
        is a right ideal of $R$.
\end{theorem}
\begin{proof}
    $\RightRegularModule{R}$ is a left module over
        $\OppositeRing{R}$ by
        \obv{right_regular_module_is_exact}.
    $\carrier[\RightRegularModule{R}]=\carrier[R]$ by
        \obv{right_regular_module_carrier}.
    For every $k\in K$ we have
        $F(k)$ is a submodule of $\RightRegularModule{R}$ over
            $\OppositeRing{R}$ by
        \obv{right_ideal_is_submodule}.
    $F\in\funs{K}{\pow{\carrier[\RightRegularModule{R}]}}$.
    Hence $\indexedIntersection{K}{F}{\carrier[R]}$
        is a submodule of $\RightRegularModule{R}$ over
            $\OppositeRing{R}$ by
        \auto{submodule_indexed_intersection}.
    Follows by \obv{right_ideal}.
\end{proof}

\begin{theorem}\label{two_sided_ideal_indexed_intersection}
    Let $R$ be a ring.
    Let $F\in\funs{K}{\pow{\carrier[R]}}$.
    Suppose for every $k\in K$ we have
        $F(k)$ is a two-sided ideal of $R$.
    Then $\indexedIntersection{K}{F}{\carrier[R]}$
        is a two-sided ideal of $R$.
\end{theorem}
\begin{proof}
    For every $k\in K$ we have $F(k)$ is a left ideal of $R$ and
        $F(k)$ is a right ideal of $R$ by \auto{two_sided_ideal}.
    Hence $\indexedIntersection{K}{F}{\carrier[R]}$
        is a left ideal of $R$ and
        $\indexedIntersection{K}{F}{\carrier[R]}$
        is a right ideal of $R$ by
        \auto{left_ideal_indexed_intersection,
            right_ideal_indexed_intersection}.
    Follows by \obv{two_sided_ideal}.
\end{proof}

\begin{corollary}\label{ideal_indexed_intersection}
    Let $R$ be a ring.
    Let $F\in\funs{K}{\pow{\carrier[R]}}$.
    Suppose for every $k\in K$ we have $F(k)$ is an ideal of $R$.
    Then $\indexedIntersection{K}{F}{\carrier[R]}$ is an ideal of $R$.
\end{corollary}
\begin{proof}
    For every $k\in K$ we have $F(k)$ is a two-sided ideal of $R$ by
        \auto{ideal}.
    Hence $\indexedIntersection{K}{F}{\carrier[R]}$
        is a two-sided ideal of $R$ by
        \obv{two_sided_ideal_indexed_intersection}.
    Follows by \obv{ideal}.
\end{proof}

\subsection{Generated ideals}

The ideal carriers of each fixed sidedness form a closure system inside the
ring carrier.  Generated ideals are the corresponding bounded hulls, so their
leastness is part of the construction rather than a later choice of
generators.

\begin{proposition}\label{left_ideal_carriers_form_closure_system}
    Let $R$ be a ring.
    Let
        $C=\{B\in\pow{\carrier[R]}\mid
            \text{$B$ is a left ideal of $R$}\}$.
    Then $C$ is a closure system on $\carrier[R]$.
\end{proposition}
\begin{proof}
    Show that $C\subseteq\pow{\carrier[R]}$.
    \begin{subproof}
        Show that for every $B\in C$ we have
            $B\in\pow{\carrier[R]}$.
        \begin{subproof}
            Follows by \auto.
        \end{subproof}
        Follows by \obv{subseteq}.
    \end{subproof}
    Show that for every $D\subseteq C$ we have
        $\boundedIntersection{\carrier[R]}{D}\in C$.
    \begin{subproof}
        Fix $D\subseteq C$.
        Show that $D\subseteq\pow{\carrier[R]}$.
        \begin{subproof}
            Show that for every $B\in D$ we have
                $B\in\pow{\carrier[R]}$.
            \begin{subproof}
                Follows by \obv{subseteq}.
            \end{subproof}
            Follows by \obv{subseteq}.
        \end{subproof}
        Hence $\identity{D}\in
            \funs{D}{\pow{\carrier[R]}}$ by
            \auto{id_elem_funs,funs_weaken_codom}.
        Show that for every $B\in D$ we have
            $(\identity{D})(B)$ is a left ideal of $R$.
        \begin{subproof}
            Fix $B\in D$.
            Hence $B\in C$ by \obv{subseteq}.
            Hence $B$ is a left ideal of $R$ by \auto.
            Follows by \obv{id_apply}.
        \end{subproof}
        Hence
            $\indexedIntersection{D}{\identity{D}}{\carrier[R]}$
            is a left ideal of $R$ by
            \obv{left_ideal_indexed_intersection}.
        $\boundedIntersection{\carrier[R]}{D}
            =\indexedIntersection{D}{\identity{D}}{\carrier[R]}$ by
            \obv{bounded_intersection_as_indexed_intersection}.
        Hence $\boundedIntersection{\carrier[R]}{D}$ is a left ideal of
            $R$.
        $\boundedIntersection{\carrier[R]}{D}
            \subseteq\carrier[R]$ by
            \obv{bounded_intersection_subseteq}.
        Hence $\boundedIntersection{\carrier[R]}{D}
            \in\pow{\carrier[R]}$ by \obv{powerset_intro}.
        Follows by \auto.
    \end{subproof}
    Follows by \obv{closure_system}.
\end{proof}

\begin{proposition}\label{right_ideal_carriers_form_closure_system}
    Let $R$ be a ring.
    Let
        $C=\{B\in\pow{\carrier[R]}\mid
            \text{$B$ is a right ideal of $R$}\}$.
    Then $C$ is a closure system on $\carrier[R]$.
\end{proposition}
\begin{proof}
    Show that $C\subseteq\pow{\carrier[R]}$.
    \begin{subproof}
        Show that for every $B\in C$ we have
            $B\in\pow{\carrier[R]}$.
        \begin{subproof}
            Follows by \auto.
        \end{subproof}
        Follows by \obv{subseteq}.
    \end{subproof}
    Show that for every $D\subseteq C$ we have
        $\boundedIntersection{\carrier[R]}{D}\in C$.
    \begin{subproof}
        Fix $D\subseteq C$.
        Show that $D\subseteq\pow{\carrier[R]}$.
        \begin{subproof}
            Show that for every $B\in D$ we have
                $B\in\pow{\carrier[R]}$.
            \begin{subproof}
                Follows by \obv{subseteq}.
            \end{subproof}
            Follows by \obv{subseteq}.
        \end{subproof}
        Hence $\identity{D}\in
            \funs{D}{\pow{\carrier[R]}}$ by
            \auto{id_elem_funs,funs_weaken_codom}.
        Show that for every $B\in D$ we have
            $(\identity{D})(B)$ is a right ideal of $R$.
        \begin{subproof}
            Fix $B\in D$.
            Hence $B\in C$ by \obv{subseteq}.
            Hence $B$ is a right ideal of $R$ by \auto.
            Follows by \obv{id_apply}.
        \end{subproof}
        Hence
            $\indexedIntersection{D}{\identity{D}}{\carrier[R]}$
            is a right ideal of $R$ by
            \obv{right_ideal_indexed_intersection}.
        $\boundedIntersection{\carrier[R]}{D}
            =\indexedIntersection{D}{\identity{D}}{\carrier[R]}$ by
            \obv{bounded_intersection_as_indexed_intersection}.
        Hence $\boundedIntersection{\carrier[R]}{D}$ is a right ideal of
            $R$.
        $\boundedIntersection{\carrier[R]}{D}
            \subseteq\carrier[R]$ by
            \obv{bounded_intersection_subseteq}.
        Hence $\boundedIntersection{\carrier[R]}{D}
            \in\pow{\carrier[R]}$ by \obv{powerset_intro}.
        Follows by \auto.
    \end{subproof}
    Follows by \obv{closure_system}.
\end{proof}

\begin{proposition}\label{two_sided_ideal_carriers_form_closure_system}
    Let $R$ be a ring.
    Let
        $C=\{B\in\pow{\carrier[R]}\mid
            \text{$B$ is a two-sided ideal of $R$}\}$.
    Then $C$ is a closure system on $\carrier[R]$.
\end{proposition}
\begin{proof}
    Let
        $L=\{B\in\pow{\carrier[R]}\mid
            \text{$B$ is a left ideal of $R$}\}$ and
        $Q=\{B\in\pow{\carrier[R]}\mid
            \text{$B$ is a right ideal of $R$}\}$.
    $L$ is a closure system on $\carrier[R]$ and
        $Q$ is a closure system on $\carrier[R]$ by
        \obv{left_ideal_carriers_form_closure_system,
            right_ideal_carriers_form_closure_system}.
    Show that $C\subseteq\pow{\carrier[R]}$.
    \begin{subproof}
        Show that for every $B\in C$ we have
            $B\in\pow{\carrier[R]}$.
        \begin{subproof}
            Follows by \auto.
        \end{subproof}
        Follows by \obv{subseteq}.
    \end{subproof}
    Show that for every $D\subseteq C$ we have
        $\boundedIntersection{\carrier[R]}{D}\in C$.
    \begin{subproof}
        Fix $D\subseteq C$.
        Show that $D\subseteq L$ and $D\subseteq Q$.
        \begin{subproof}
            Show that for every $B\in D$ we have $B\in L$ and $B\in Q$.
            \begin{subproof}
                Fix $B\in D$.
                Hence $B\in C$ by \obv{subseteq}.
                Hence $B\in\pow{\carrier[R]}$ and
                    $B$ is a two-sided ideal of $R$ by \auto.
                Hence $B$ is a left ideal of $R$ and
                    $B$ is a right ideal of $R$ by
                    \obv{two_sided_ideal}.
            \end{subproof}
            Follows by \auto{subseteq}.
        \end{subproof}
        Hence $\boundedIntersection{\carrier[R]}{D}\in L$ and
            $\boundedIntersection{\carrier[R]}{D}\in Q$ by
            \obv{closure_system_intersection}.
        Hence $\boundedIntersection{\carrier[R]}{D}$ is a left ideal of
            $R$ and $\boundedIntersection{\carrier[R]}{D}$ is a right
            ideal of $R$.
        Hence $\boundedIntersection{\carrier[R]}{D}$ is a two-sided ideal
            of $R$ by \obv{two_sided_ideal}.
        $\boundedIntersection{\carrier[R]}{D}
            \subseteq\carrier[R]$ by
            \obv{bounded_intersection_subseteq}.
        Hence $\boundedIntersection{\carrier[R]}{D}
            \in\pow{\carrier[R]}$ by \obv{powerset_intro}.
        Follows by \auto.
    \end{subproof}
    Follows by \obv{closure_system}.
\end{proof}

\begin{definition}\label{generated_left_ideal}
    Let $R$ be a ring.
    $\generatedLeftIdeal{R}{A}=
        \closureHull{\carrier[R]}
            {\{B\in\pow{\carrier[R]}\mid
                \text{$B$ is a left ideal of $R$}\}}
            {A}$.
\end{definition}

\begin{definition}\label{generated_right_ideal}
    Let $R$ be a ring.
    $\generatedRightIdeal{R}{A}=
        \closureHull{\carrier[R]}
            {\{B\in\pow{\carrier[R]}\mid
                \text{$B$ is a right ideal of $R$}\}}
            {A}$.
\end{definition}

\begin{definition}\label{generated_two_sided_ideal}
    Let $R$ be a ring.
    $\generatedTwoSidedIdeal{R}{A}=
        \closureHull{\carrier[R]}
            {\{B\in\pow{\carrier[R]}\mid
                \text{$B$ is a two-sided ideal of $R$}\}}
            {A}$.
\end{definition}

\begin{definition}\label{generated_ideal}
    Let $R$ be a ring.
    $\generatedIdeal{R}{A}=\generatedTwoSidedIdeal{R}{A}$.
\end{definition}

\subsubsection{Left-generated ideals}

\begin{proposition}\label{generated_left_ideal_subseteq_carrier}
    Let $R$ be a ring.
    Let $A$ be a subset of $\carrier[R]$.
    Then $\generatedLeftIdeal{R}{A}\subseteq\carrier[R]$.
\end{proposition}
\begin{proof}
    Follows by \auto{generated_left_ideal,closure_hull_subseteq}.
\end{proof}

\begin{proposition}\label{generated_left_ideal_contains_generators}
    Let $R$ be a ring.
    Let $A$ be a subset of $\carrier[R]$.
    Then $A\subseteq\generatedLeftIdeal{R}{A}$.
\end{proposition}
\begin{proof}
    Let
        $C=\{B\in\pow{\carrier[R]}\mid
            \text{$B$ is a left ideal of $R$}\}$.
    $C$ is a closure system on $\carrier[R]$ by
        \obv{left_ideal_carriers_form_closure_system}.
    $A\subseteq\closureHull{\carrier[R]}{C}{A}$ by
        \obv{closure_hull_extensive}.
    Follows by \auto{generated_left_ideal}.
\end{proof}

\begin{theorem}\label{generated_left_ideal_is_left_ideal}
    Let $R$ be a ring.
    Let $A$ be a subset of $\carrier[R]$.
    Then $\generatedLeftIdeal{R}{A}$ is a left ideal of $R$.
\end{theorem}
\begin{proof}
    Let
        $C=\{B\in\pow{\carrier[R]}\mid
            \text{$B$ is a left ideal of $R$}\}$.
    $C$ is a closure system on $\carrier[R]$ by
        \obv{left_ideal_carriers_form_closure_system}.
    $\closureHull{\carrier[R]}{C}{A}\in C$ by
        \obv{closure_hull_closed}.
    Follows by \auto{generated_left_ideal}.
\end{proof}

\begin{theorem}\label{generated_left_ideal_least}
    Let $R$ be a ring.
    Let $A$ be a subset of $\carrier[R]$.
    Suppose $I$ is a left ideal of $R$ and $A\subseteq I$.
    Then $\generatedLeftIdeal{R}{A}\subseteq I$.
\end{theorem}
\begin{proof}
    Let
        $C=\{B\in\pow{\carrier[R]}\mid
            \text{$B$ is a left ideal of $R$}\}$.
    $C$ is a closure system on $\carrier[R]$ by
        \obv{left_ideal_carriers_form_closure_system}.
    $I\subseteq\carrier[R]$ by \obv{left_ideal_subseteq_carrier}.
    Hence $I\in\pow{\carrier[R]}$ by \obv{powerset_intro}.
    Hence $I\in C$.
    $\closureHull{\carrier[R]}{C}{A}\subseteq I$ by
        \obv{closure_hull_least}.
    Follows by \auto{generated_left_ideal}.
\end{proof}

\begin{theorem}\label{generated_left_ideal_universal_property}
    Let $R$ be a ring.
    Let $A$ be a subset of $\carrier[R]$.
    Suppose $I$ is a left ideal of $R$.
    Then $\generatedLeftIdeal{R}{A}\subseteq I$ iff $A\subseteq I$.
\end{theorem}
\begin{proof}
    $A\subseteq\generatedLeftIdeal{R}{A}$ by
        \obv{generated_left_ideal_contains_generators}.
    Follows by
        \auto{generated_left_ideal_least,subseteq_transitive}.
\end{proof}

\begin{proposition}\label{generated_left_ideal_monotone}
    Let $R$ be a ring.
    Let $A,B\subseteq\carrier[R]$.
    Suppose $A\subseteq B$.
    Then $\generatedLeftIdeal{R}{A}\subseteq
        \generatedLeftIdeal{R}{B}$.
\end{proposition}
\begin{proof}
    $\generatedLeftIdeal{R}{B}$ is a left ideal of $R$ by
        \obv{generated_left_ideal_is_left_ideal}.
    $B\subseteq\generatedLeftIdeal{R}{B}$ by
        \obv{generated_left_ideal_contains_generators}.
    Hence $A\subseteq\generatedLeftIdeal{R}{B}$ by
        \obv{subseteq_transitive}.
    Follows by \obv{generated_left_ideal_least}.
\end{proof}

\begin{corollary}\label{generated_left_ideal_of_left_ideal}
    Let $R$ be a ring.
    Suppose $I$ is a left ideal of $R$.
    Then $\generatedLeftIdeal{R}{I}=I$.
\end{corollary}
\begin{proof}
    $I\subseteq\carrier[R]$ by \obv{left_ideal_subseteq_carrier}.
    $I\subseteq\generatedLeftIdeal{R}{I}$ by
        \obv{generated_left_ideal_contains_generators}.
    $\generatedLeftIdeal{R}{I}\subseteq I$ by
        \obv{generated_left_ideal_least,subseteq_refl}.
    Follows by \obv{subseteq_antisymmetric}.
\end{proof}

\begin{theorem}\label{principal_left_ideal_properties}
    Let $R$ be a ring and $x\in\carrier[R]$.
    Then $\generatedLeftIdeal{R}{\{x\}}$ is a left ideal of $R$ and
        $x\in\generatedLeftIdeal{R}{\{x\}}$ and
        for every left ideal $I$ of $R$ we have
            $\generatedLeftIdeal{R}{\{x\}}\subseteq I$ iff $x\in I$.
\end{theorem}
\begin{proof}
    $\{x\}\subseteq\carrier[R]$ by \obv{singleton_subset_intro}.
    Hence $\generatedLeftIdeal{R}{\{x\}}$ is a left ideal of $R$ and
        $\{x\}\subseteq\generatedLeftIdeal{R}{\{x\}}$ by
        \obv{generated_left_ideal_is_left_ideal,
            generated_left_ideal_contains_generators}.
    Hence $x\in\generatedLeftIdeal{R}{\{x\}}$ by
        \obv{singleton_subset_elim}.
    For every left ideal $I$ of $R$ we have
        $\generatedLeftIdeal{R}{\{x\}}\subseteq I$ iff
        $\{x\}\subseteq I$ by
        \obv{generated_left_ideal_universal_property}.
    For every set $I$ we have $\{x\}\subseteq I$ iff $x\in I$ by
        \auto{singleton_subset_intro,singleton_subset_elim}.
    Follows by \auto.
\end{proof}

\subsubsection{Right-generated ideals}

\begin{proposition}\label{generated_right_ideal_subseteq_carrier}
    Let $R$ be a ring.
    Let $A$ be a subset of $\carrier[R]$.
    Then $\generatedRightIdeal{R}{A}\subseteq\carrier[R]$.
\end{proposition}
\begin{proof}
    Follows by \auto{generated_right_ideal,closure_hull_subseteq}.
\end{proof}

\begin{proposition}\label{generated_right_ideal_contains_generators}
    Let $R$ be a ring.
    Let $A$ be a subset of $\carrier[R]$.
    Then $A\subseteq\generatedRightIdeal{R}{A}$.
\end{proposition}
\begin{proof}
    Let
        $C=\{B\in\pow{\carrier[R]}\mid
            \text{$B$ is a right ideal of $R$}\}$.
    $C$ is a closure system on $\carrier[R]$ by
        \obv{right_ideal_carriers_form_closure_system}.
    $A\subseteq\closureHull{\carrier[R]}{C}{A}$ by
        \obv{closure_hull_extensive}.
    Follows by \auto{generated_right_ideal}.
\end{proof}

\begin{theorem}\label{generated_right_ideal_is_right_ideal}
    Let $R$ be a ring.
    Let $A$ be a subset of $\carrier[R]$.
    Then $\generatedRightIdeal{R}{A}$ is a right ideal of $R$.
\end{theorem}
\begin{proof}
    Let
        $C=\{B\in\pow{\carrier[R]}\mid
            \text{$B$ is a right ideal of $R$}\}$.
    $C$ is a closure system on $\carrier[R]$ by
        \obv{right_ideal_carriers_form_closure_system}.
    $\closureHull{\carrier[R]}{C}{A}\in C$ by
        \obv{closure_hull_closed}.
    Follows by \auto{generated_right_ideal}.
\end{proof}

\begin{theorem}\label{generated_right_ideal_least}
    Let $R$ be a ring.
    Let $A$ be a subset of $\carrier[R]$.
    Suppose $I$ is a right ideal of $R$ and $A\subseteq I$.
    Then $\generatedRightIdeal{R}{A}\subseteq I$.
\end{theorem}
\begin{proof}
    Let
        $C=\{B\in\pow{\carrier[R]}\mid
            \text{$B$ is a right ideal of $R$}\}$.
    $C$ is a closure system on $\carrier[R]$ by
        \obv{right_ideal_carriers_form_closure_system}.
    $I\subseteq\carrier[R]$ by \obv{right_ideal_subseteq_carrier}.
    Hence $I\in\pow{\carrier[R]}$ by \obv{powerset_intro}.
    Hence $I\in C$.
    $\closureHull{\carrier[R]}{C}{A}\subseteq I$ by
        \obv{closure_hull_least}.
    Follows by \auto{generated_right_ideal}.
\end{proof}

\begin{theorem}\label{generated_right_ideal_universal_property}
    Let $R$ be a ring.
    Let $A$ be a subset of $\carrier[R]$.
    Suppose $I$ is a right ideal of $R$.
    Then $\generatedRightIdeal{R}{A}\subseteq I$ iff $A\subseteq I$.
\end{theorem}
\begin{proof}
    $A\subseteq\generatedRightIdeal{R}{A}$ by
        \obv{generated_right_ideal_contains_generators}.
    Follows by
        \auto{generated_right_ideal_least,subseteq_transitive}.
\end{proof}

\begin{proposition}\label{generated_right_ideal_monotone}
    Let $R$ be a ring.
    Let $A,B\subseteq\carrier[R]$.
    Suppose $A\subseteq B$.
    Then $\generatedRightIdeal{R}{A}\subseteq
        \generatedRightIdeal{R}{B}$.
\end{proposition}
\begin{proof}
    $\generatedRightIdeal{R}{B}$ is a right ideal of $R$ by
        \obv{generated_right_ideal_is_right_ideal}.
    $B\subseteq\generatedRightIdeal{R}{B}$ by
        \obv{generated_right_ideal_contains_generators}.
    Hence $A\subseteq\generatedRightIdeal{R}{B}$ by
        \obv{subseteq_transitive}.
    Follows by \obv{generated_right_ideal_least}.
\end{proof}

\begin{corollary}\label{generated_right_ideal_of_right_ideal}
    Let $R$ be a ring.
    Suppose $I$ is a right ideal of $R$.
    Then $\generatedRightIdeal{R}{I}=I$.
\end{corollary}
\begin{proof}
    $I\subseteq\carrier[R]$ by \obv{right_ideal_subseteq_carrier}.
    $I\subseteq\generatedRightIdeal{R}{I}$ by
        \obv{generated_right_ideal_contains_generators}.
    $\generatedRightIdeal{R}{I}\subseteq I$ by
        \obv{generated_right_ideal_least,subseteq_refl}.
    Follows by \obv{subseteq_antisymmetric}.
\end{proof}

\begin{theorem}\label{principal_right_ideal_properties}
    Let $R$ be a ring and $x\in\carrier[R]$.
    Then $\generatedRightIdeal{R}{\{x\}}$ is a right ideal of $R$ and
        $x\in\generatedRightIdeal{R}{\{x\}}$ and
        for every right ideal $I$ of $R$ we have
            $\generatedRightIdeal{R}{\{x\}}\subseteq I$ iff $x\in I$.
\end{theorem}
\begin{proof}
    $\{x\}\subseteq\carrier[R]$ by \obv{singleton_subset_intro}.
    Hence $\generatedRightIdeal{R}{\{x\}}$ is a right ideal of $R$ and
        $\{x\}\subseteq\generatedRightIdeal{R}{\{x\}}$ by
        \obv{generated_right_ideal_is_right_ideal,
            generated_right_ideal_contains_generators}.
    Hence $x\in\generatedRightIdeal{R}{\{x\}}$ by
        \obv{singleton_subset_elim}.
    For every right ideal $I$ of $R$ we have
        $\generatedRightIdeal{R}{\{x\}}\subseteq I$ iff
        $\{x\}\subseteq I$ by
        \obv{generated_right_ideal_universal_property}.
    For every set $I$ we have $\{x\}\subseteq I$ iff $x\in I$ by
        \auto{singleton_subset_intro,singleton_subset_elim}.
    Follows by \auto.
\end{proof}

\subsubsection{Two-sided generated ideals}

\begin{proposition}\label{generated_two_sided_ideal_subseteq_carrier}
    Let $R$ be a ring.
    Let $A$ be a subset of $\carrier[R]$.
    Then $\generatedTwoSidedIdeal{R}{A}\subseteq\carrier[R]$.
\end{proposition}
\begin{proof}
    Follows by
        \auto{generated_two_sided_ideal,closure_hull_subseteq}.
\end{proof}

\begin{proposition}\label{generated_two_sided_ideal_contains_generators}
    Let $R$ be a ring.
    Let $A$ be a subset of $\carrier[R]$.
    Then $A\subseteq\generatedTwoSidedIdeal{R}{A}$.
\end{proposition}
\begin{proof}
    Let
        $C=\{B\in\pow{\carrier[R]}\mid
            \text{$B$ is a two-sided ideal of $R$}\}$.
    $C$ is a closure system on $\carrier[R]$ by
        \obv{two_sided_ideal_carriers_form_closure_system}.
    $A\subseteq\closureHull{\carrier[R]}{C}{A}$ by
        \obv{closure_hull_extensive}.
    Follows by \auto{generated_two_sided_ideal}.
\end{proof}

\begin{theorem}\label{generated_two_sided_ideal_is_two_sided_ideal}
    Let $R$ be a ring.
    Let $A$ be a subset of $\carrier[R]$.
    Then $\generatedTwoSidedIdeal{R}{A}$ is a two-sided ideal of $R$.
\end{theorem}
\begin{proof}
    Let
        $C=\{B\in\pow{\carrier[R]}\mid
            \text{$B$ is a two-sided ideal of $R$}\}$.
    $C$ is a closure system on $\carrier[R]$ by
        \obv{two_sided_ideal_carriers_form_closure_system}.
    $\closureHull{\carrier[R]}{C}{A}\in C$ by
        \obv{closure_hull_closed}.
    Follows by \auto{generated_two_sided_ideal}.
\end{proof}

\begin{theorem}\label{generated_two_sided_ideal_least}
    Let $R$ be a ring.
    Let $A$ be a subset of $\carrier[R]$.
    Suppose $I$ is a two-sided ideal of $R$ and $A\subseteq I$.
    Then $\generatedTwoSidedIdeal{R}{A}\subseteq I$.
\end{theorem}
\begin{proof}
    Let
        $C=\{B\in\pow{\carrier[R]}\mid
            \text{$B$ is a two-sided ideal of $R$}\}$.
    $C$ is a closure system on $\carrier[R]$ by
        \obv{two_sided_ideal_carriers_form_closure_system}.
    $I\subseteq\carrier[R]$ by
        \obv{two_sided_ideal_subseteq_carrier}.
    Hence $I\in\pow{\carrier[R]}$ by \obv{powerset_intro}.
    Hence $I\in C$.
    $\closureHull{\carrier[R]}{C}{A}\subseteq I$ by
        \obv{closure_hull_least}.
    Follows by \auto{generated_two_sided_ideal}.
\end{proof}

\begin{theorem}\label{generated_two_sided_ideal_universal_property}
    Let $R$ be a ring.
    Let $A$ be a subset of $\carrier[R]$.
    Suppose $I$ is a two-sided ideal of $R$.
    Then $\generatedTwoSidedIdeal{R}{A}\subseteq I$ iff $A\subseteq I$.
\end{theorem}
\begin{proof}
    $A\subseteq\generatedTwoSidedIdeal{R}{A}$ by
        \obv{generated_two_sided_ideal_contains_generators}.
    Follows by
        \auto{generated_two_sided_ideal_least,subseteq_transitive}.
\end{proof}

\begin{proposition}\label{generated_two_sided_ideal_monotone}
    Let $R$ be a ring.
    Let $A,B\subseteq\carrier[R]$.
    Suppose $A\subseteq B$.
    Then $\generatedTwoSidedIdeal{R}{A}\subseteq
        \generatedTwoSidedIdeal{R}{B}$.
\end{proposition}
\begin{proof}
    $\generatedTwoSidedIdeal{R}{B}$ is a two-sided ideal of $R$ by
        \obv{generated_two_sided_ideal_is_two_sided_ideal}.
    $B\subseteq\generatedTwoSidedIdeal{R}{B}$ by
        \obv{generated_two_sided_ideal_contains_generators}.
    Hence $A\subseteq\generatedTwoSidedIdeal{R}{B}$ by
        \obv{subseteq_transitive}.
    Follows by \obv{generated_two_sided_ideal_least}.
\end{proof}

\begin{corollary}\label{generated_two_sided_ideal_of_two_sided_ideal}
    Let $R$ be a ring.
    Suppose $I$ is a two-sided ideal of $R$.
    Then $\generatedTwoSidedIdeal{R}{I}=I$.
\end{corollary}
\begin{proof}
    $I\subseteq\carrier[R]$ by
        \obv{two_sided_ideal_subseteq_carrier}.
    $I\subseteq\generatedTwoSidedIdeal{R}{I}$ by
        \obv{generated_two_sided_ideal_contains_generators}.
    $\generatedTwoSidedIdeal{R}{I}\subseteq I$ by
        \obv{generated_two_sided_ideal_least,subseteq_refl}.
    Follows by \obv{subseteq_antisymmetric}.
\end{proof}

\begin{theorem}\label{principal_two_sided_ideal_properties}
    Let $R$ be a ring and $x\in\carrier[R]$.
    Then $\generatedTwoSidedIdeal{R}{\{x\}}$ is a two-sided ideal of $R$
        and $x\in\generatedTwoSidedIdeal{R}{\{x\}}$ and
        for every two-sided ideal $I$ of $R$ we have
            $\generatedTwoSidedIdeal{R}{\{x\}}\subseteq I$ iff $x\in I$.
\end{theorem}
\begin{proof}
    $\{x\}\subseteq\carrier[R]$ by \obv{singleton_subset_intro}.
    Hence $\generatedTwoSidedIdeal{R}{\{x\}}$ is a two-sided ideal of $R$
        and $\{x\}\subseteq\generatedTwoSidedIdeal{R}{\{x\}}$ by
        \obv{generated_two_sided_ideal_is_two_sided_ideal,
            generated_two_sided_ideal_contains_generators}.
    Hence $x\in\generatedTwoSidedIdeal{R}{\{x\}}$ by
        \obv{singleton_subset_elim}.
    For every two-sided ideal $I$ of $R$ we have
        $\generatedTwoSidedIdeal{R}{\{x\}}\subseteq I$ iff
        $\{x\}\subseteq I$ by
        \obv{generated_two_sided_ideal_universal_property}.
    For every set $I$ we have $\{x\}\subseteq I$ iff $x\in I$ by
        \auto{singleton_subset_intro,singleton_subset_elim}.
    Follows by \auto.
\end{proof}

\subsubsection{Generated ideals}

\begin{proposition}\label{generated_ideal_subseteq_carrier}
    Let $R$ be a ring.
    Let $A$ be a subset of $\carrier[R]$.
    Then $\generatedIdeal{R}{A}\subseteq\carrier[R]$.
\end{proposition}
\begin{proof}
    Follows by
        \auto{generated_ideal,
            generated_two_sided_ideal_subseteq_carrier}.
\end{proof}

\begin{proposition}\label{generated_ideal_contains_generators}
    Let $R$ be a ring.
    Let $A$ be a subset of $\carrier[R]$.
    Then $A\subseteq\generatedIdeal{R}{A}$.
\end{proposition}
\begin{proof}
    Follows by
        \auto{generated_ideal,
            generated_two_sided_ideal_contains_generators}.
\end{proof}

\begin{theorem}\label{generated_ideal_is_ideal}
    Let $R$ be a ring.
    Let $A$ be a subset of $\carrier[R]$.
    Then $\generatedIdeal{R}{A}$ is an ideal of $R$.
\end{theorem}
\begin{proof}
    $\generatedTwoSidedIdeal{R}{A}$ is a two-sided ideal of $R$ by
        \obv{generated_two_sided_ideal_is_two_sided_ideal}.
    Follows by \auto{generated_ideal,ideal}.
\end{proof}

\begin{theorem}\label{generated_ideal_least}
    Let $R$ be a ring.
    Let $A$ be a subset of $\carrier[R]$.
    Suppose $I$ is an ideal of $R$ and $A\subseteq I$.
    Then $\generatedIdeal{R}{A}\subseteq I$.
\end{theorem}
\begin{proof}
    $I$ is a two-sided ideal of $R$ by \obv{ideal}.
    Follows by
        \auto{generated_ideal,generated_two_sided_ideal_least}.
\end{proof}

\begin{theorem}\label{generated_ideal_universal_property}
    Let $R$ be a ring.
    Let $A$ be a subset of $\carrier[R]$.
    Suppose $I$ is an ideal of $R$.
    Then $\generatedIdeal{R}{A}\subseteq I$ iff $A\subseteq I$.
\end{theorem}
\begin{proof}
    $A\subseteq\generatedIdeal{R}{A}$ by
        \obv{generated_ideal_contains_generators}.
    Follows by \auto{generated_ideal_least,subseteq_transitive}.
\end{proof}

\begin{proposition}\label{generated_ideal_monotone}
    Let $R$ be a ring.
    Let $A,B\subseteq\carrier[R]$.
    Suppose $A\subseteq B$.
    Then $\generatedIdeal{R}{A}\subseteq\generatedIdeal{R}{B}$.
\end{proposition}
\begin{proof}
    Follows by
        \auto{generated_ideal,generated_two_sided_ideal_monotone}.
\end{proof}

\begin{theorem}\label{generated_ideal_eq_generated_two_sided_ideal}
    Let $R$ be a ring.
    Let $A$ be a subset of $\carrier[R]$.
    Then
        $\generatedIdeal{R}{A}=\generatedTwoSidedIdeal{R}{A}$.
\end{theorem}
\begin{proof}
    Follows by \obv{generated_ideal}.
\end{proof}

\begin{theorem}\label{commutative_generated_ideal_eq_generated_left_ideal}
    Let $R$ be a commutative ring.
    Let $A$ be a subset of $\carrier[R]$.
    Then $\generatedIdeal{R}{A}=\generatedLeftIdeal{R}{A}$.
\end{theorem}
\begin{proof}
    $R$ is a ring by \obv{commutative_ring_is_ring}.
    $\generatedIdeal{R}{A}$ is an ideal of $R$ by
        \obv{generated_ideal_is_ideal}.
    Hence $\generatedIdeal{R}{A}$ is a left ideal of $R$ by
        \auto{commutative_ideal_iff_left_ideal}.
    $A\subseteq\generatedIdeal{R}{A}$ by
        \obv{generated_ideal_contains_generators}.
    Hence $\generatedLeftIdeal{R}{A}\subseteq
        \generatedIdeal{R}{A}$ by \obv{generated_left_ideal_least}.
    $\generatedLeftIdeal{R}{A}$ is a left ideal of $R$ by
        \obv{generated_left_ideal_is_left_ideal}.
    Hence $\generatedLeftIdeal{R}{A}$ is an ideal of $R$ by
        \auto{commutative_ideal_iff_left_ideal}.
    $A\subseteq\generatedLeftIdeal{R}{A}$ by
        \obv{generated_left_ideal_contains_generators}.
    Hence $\generatedIdeal{R}{A}\subseteq
        \generatedLeftIdeal{R}{A}$ by \obv{generated_ideal_least}.
    Follows by \obv{subseteq_antisymmetric}.
\end{proof}

\begin{theorem}\label{commutative_generated_ideal_eq_generated_right_ideal}
    Let $R$ be a commutative ring.
    Let $A$ be a subset of $\carrier[R]$.
    Then $\generatedIdeal{R}{A}=\generatedRightIdeal{R}{A}$.
\end{theorem}
\begin{proof}
    $R$ is a ring by \obv{commutative_ring_is_ring}.
    $\generatedIdeal{R}{A}$ is an ideal of $R$ by
        \obv{generated_ideal_is_ideal}.
    Hence $\generatedIdeal{R}{A}$ is a right ideal of $R$ by
        \auto{commutative_ideal_iff_right_ideal}.
    $A\subseteq\generatedIdeal{R}{A}$ by
        \obv{generated_ideal_contains_generators}.
    Hence $\generatedRightIdeal{R}{A}\subseteq
        \generatedIdeal{R}{A}$ by \obv{generated_right_ideal_least}.
    $\generatedRightIdeal{R}{A}$ is a right ideal of $R$ by
        \obv{generated_right_ideal_is_right_ideal}.
    Hence $\generatedRightIdeal{R}{A}$ is an ideal of $R$ by
        \auto{commutative_ideal_iff_right_ideal}.
    $A\subseteq\generatedRightIdeal{R}{A}$ by
        \obv{generated_right_ideal_contains_generators}.
    Hence $\generatedIdeal{R}{A}\subseteq
        \generatedRightIdeal{R}{A}$ by \obv{generated_ideal_least}.
    Follows by \obv{subseteq_antisymmetric}.
\end{proof}

\begin{corollary}\label{generated_ideal_of_ideal}
    Let $R$ be a ring.
    Suppose $I$ is an ideal of $R$.
    Then $\generatedIdeal{R}{I}=I$.
\end{corollary}
\begin{proof}
    $I\subseteq\carrier[R]$ by
        \obv{ideal,two_sided_ideal_subseteq_carrier}.
    $I\subseteq\generatedIdeal{R}{I}$ by
        \obv{generated_ideal_contains_generators}.
    $\generatedIdeal{R}{I}\subseteq I$ by
        \obv{generated_ideal_least,subseteq_refl}.
    Follows by \obv{subseteq_antisymmetric}.
\end{proof}

\begin{theorem}\label{principal_ideal_properties}
    Let $R$ be a ring and $x\in\carrier[R]$.
    Then $\generatedIdeal{R}{\{x\}}$ is an ideal of $R$ and
        $x\in\generatedIdeal{R}{\{x\}}$ and
        for every ideal $I$ of $R$ we have
            $\generatedIdeal{R}{\{x\}}\subseteq I$ iff $x\in I$.
\end{theorem}
\begin{proof}
    $\{x\}\subseteq\carrier[R]$ by \obv{singleton_subset_intro}.
    Hence $\generatedIdeal{R}{\{x\}}$ is an ideal of $R$ and
        $\{x\}\subseteq\generatedIdeal{R}{\{x\}}$ by
        \obv{generated_ideal_is_ideal,
            generated_ideal_contains_generators}.
    Hence $x\in\generatedIdeal{R}{\{x\}}$ by
        \obv{singleton_subset_elim}.
    For every ideal $I$ of $R$ we have
        $\generatedIdeal{R}{\{x\}}\subseteq I$ iff
        $\{x\}\subseteq I$ by
        \obv{generated_ideal_universal_property}.
    For every set $I$ we have $\{x\}\subseteq I$ iff $x\in I$ by
        \auto{singleton_subset_intro,singleton_subset_elim}.
    Follows by \auto.
\end{proof}

\subsubsection{Principal ideals}

\begin{definition}\label{principal_ideal}
    Let $R$ be a ring.
    $\principalIdeal{R}{a}=\generatedIdeal{R}{\{a\}}$.
\end{definition}

\begin{theorem}\label{principal_ideal_is_ideal}
    Let $R$ be a ring and $a\in\carrier[R]$.
    Then $\principalIdeal{R}{a}$ is an ideal of $R$.
\end{theorem}
\begin{proof}
    Follows by
        \auto{principal_ideal,principal_ideal_properties}.
\end{proof}

\begin{proposition}\label{principal_ideal_contains_generator}
    Let $R$ be a ring and $a\in\carrier[R]$.
    Then $a\in\principalIdeal{R}{a}$.
\end{proposition}
\begin{proof}
    Follows by
        \auto{principal_ideal,principal_ideal_properties}.
\end{proof}

\begin{theorem}\label{principal_ideal_universal_property}
    Let $R$ be a ring and $a\in\carrier[R]$.
    Suppose $I$ is an ideal of $R$.
    Then $\principalIdeal{R}{a}\subseteq I$ iff $a\in I$.
\end{theorem}
\begin{proof}
    Follows by
        \auto{principal_ideal,principal_ideal_properties}.
\end{proof}

\begin{theorem}\label{principal_ideal_zero}
    Let $R$ be a ring.
    Then $\principalIdeal{R}{\ringZero[R]}=\{\ringZero[R]\}$.
\end{theorem}
\begin{proof}
    $\{\ringZero[R]\}$ is an ideal of $R$ by
        \obv{zero_is_ideal}.
    Hence $\generatedIdeal{R}{\{\ringZero[R]\}}
        =\{\ringZero[R]\}$ by \obv{generated_ideal_of_ideal}.
    Follows by \obv{principal_ideal}.
\end{proof}

\begin{theorem}\label{principal_ideal_one}
    Let $R$ be a ring.
    Then $\principalIdeal{R}{\ringOne[R]}=\carrier[R]$.
\end{theorem}
\begin{proof}
    $\ringOne[R]\in\carrier[R]$ by \obv{ring_one_type}.
    Hence $\principalIdeal{R}{\ringOne[R]}$ is an ideal of $R$ and
        $\ringOne[R]\in\principalIdeal{R}{\ringOne[R]}$ by
        \obv{principal_ideal_is_ideal,
            principal_ideal_contains_generator}.
    Hence $\principalIdeal{R}{\ringOne[R]}$ is a left ideal of $R$ by
        \auto{ideal,two_sided_ideal}.
    Follows by \obv{left_ideal_contains_one_iff_full}.
\end{proof}

\begin{theorem}\label{principal_ideal_mul_unit_left}
    Let $R$ be a ring and $a\in\carrier[R]$.
    Suppose $u$ is a unit of $\MultiplicativeMonoid{R}$.
    Then
        $\principalIdeal{R}{\ringMul[R](u,a)}
            =\principalIdeal{R}{a}$.
\end{theorem}
\begin{proof}
    $u\in\carrier[\MultiplicativeMonoid{R}]$ by
        \obv{monoid_unit_in_carrier}.
    Take $v\in\carrier[\MultiplicativeMonoid{R}]$ such that
        $\mul[\MultiplicativeMonoid{R}](u,v)
            =\neutral[\MultiplicativeMonoid{R}]$ and
        $\mul[\MultiplicativeMonoid{R}](v,u)
            =\neutral[\MultiplicativeMonoid{R}]$ by
        \obv{monoid_unit_has_inverse}.
    Hence $u,v\in\carrier[R]$ and
        $\ringMul[R](u,v)=\ringOne[R]$ and
        $\ringMul[R](v,u)=\ringOne[R]$ by
        \auto{multiplicative_monoid_fields}.
    $\ringMul[R](u,a)\in\carrier[R]$ by \obv{ring_mul_type}.

    Let $I=\principalIdeal{R}{a}$ and
        $J=\principalIdeal{R}{\ringMul[R](u,a)}$.
    $I$ and $J$ are ideals of $R$ by
        \auto{principal_ideal_is_ideal}.
    Hence $I$ and $J$ are left ideals of $R$ by
        \auto{ideal,two_sided_ideal}.
    $a\in I$ and $\ringMul[R](u,a)\in J$ by
        \auto{principal_ideal_contains_generator}.
    Hence $\ringMul[R](u,a)\in I$ by
        \auto{left_ideal_expanded_iff}.
    Hence $J\subseteq I$ by
        \auto{principal_ideal_universal_property}.

    $\ringMul[R](v,\ringMul[R](u,a))\in J$ by
        \auto{left_ideal_expanded_iff}.
    $\ringMul[R](v,\ringMul[R](u,a))=a$ by
        \auto{ring_mul_associative,ring_mul_one_left}.
    Hence $a\in J$.
    Hence $I\subseteq J$ by
        \auto{principal_ideal_universal_property}.
    Hence $I=J$ by \obv{subseteq_antisymmetric}.
    Follows by \auto.
\end{proof}

\begin{theorem}\label{principal_ideal_mul_unit_right}
    Let $R$ be a ring and $a\in\carrier[R]$.
    Suppose $u$ is a unit of $\MultiplicativeMonoid{R}$.
    Then
        $\principalIdeal{R}{\ringMul[R](a,u)}
            =\principalIdeal{R}{a}$.
\end{theorem}
\begin{proof}
    $u\in\carrier[\MultiplicativeMonoid{R}]$ by
        \obv{monoid_unit_in_carrier}.
    Take $v\in\carrier[\MultiplicativeMonoid{R}]$ such that
        $\mul[\MultiplicativeMonoid{R}](u,v)
            =\neutral[\MultiplicativeMonoid{R}]$ and
        $\mul[\MultiplicativeMonoid{R}](v,u)
            =\neutral[\MultiplicativeMonoid{R}]$ by
        \obv{monoid_unit_has_inverse}.
    Hence $u,v\in\carrier[R]$ and
        $\ringMul[R](u,v)=\ringOne[R]$ and
        $\ringMul[R](v,u)=\ringOne[R]$ by
        \auto{multiplicative_monoid_fields}.
    $\ringMul[R](a,u)\in\carrier[R]$ by \obv{ring_mul_type}.

    Let $I=\principalIdeal{R}{a}$ and
        $J=\principalIdeal{R}{\ringMul[R](a,u)}$.
    $I$ and $J$ are ideals of $R$ by
        \auto{principal_ideal_is_ideal}.
    Hence $I$ and $J$ are right ideals of $R$ by
        \auto{ideal,two_sided_ideal}.
    $a\in I$ and $\ringMul[R](a,u)\in J$ by
        \auto{principal_ideal_contains_generator}.
    Hence $\ringMul[R](a,u)\in I$ by
        \auto{right_ideal_expanded_iff}.
    Hence $J\subseteq I$ by
        \auto{principal_ideal_universal_property}.

    $\ringMul[R](\ringMul[R](a,u),v)\in J$ by
        \auto{right_ideal_expanded_iff}.
    $\ringMul[R](\ringMul[R](a,u),v)=a$ by
        \auto{ring_mul_associative,ring_mul_one_right}.
    Hence $a\in J$.
    Hence $I\subseteq J$ by
        \auto{principal_ideal_universal_property}.
    Hence $I=J$ by \obv{subseteq_antisymmetric}.
    Follows by \auto.
\end{proof}

\begin{theorem}\label{principal_ideal_eq_multiple_carrier}
    Let $R$ be a commutative ring and $a\in\carrier[R]$.
    Then
        $\principalIdeal{R}{a}
            =\{\ringMul[R](a,c)\mid c\in\carrier[R]\}$.
\end{theorem}
\begin{proof}
    $R$ is a ring by \obv{commutative_ring_is_ring}.
    Let $G=\AdditiveGroup{R}$ and
        $B=\{\ringMul[R](a,c)\mid c\in\carrier[R]\}$.
    $G$ is a group by \obv{additive_group_is_group}.
    $\carrier[G]=\carrier[R]$ and
        $\mul[G]=\ringAdd[R]$ and
        $\neutral[G]=\ringZero[R]$ and
        $\groupInv[G]=\ringNeg[R]$ by
        \auto{additive_group_fields}.

    Show that $B\subseteq\carrier[G]$.
    \begin{subproof}
        Show that every element of $B$ is an element of $\carrier[G]$.
        \begin{subproof}
            Fix $x\in B$.
            Take $c\in\carrier[R]$ such that
                $x=\ringMul[R](a,c)$ by \auto.
            $\ringMul[R](a,c)\in\carrier[R]$ by \obv{ring_mul_type}.
            Follows by \auto.
        \end{subproof}
        Follows by \obv{subseteq}.
    \end{subproof}

    $\ringZero[R]=\ringMul[R](a,\ringZero[R])$ by
        \obv{ring_mul_zero_right}.
    $\ringZero[R]\in\carrier[R]$ by \obv{ring_zero_type}.
    Hence $\neutral[G]\in B$ by \auto.

    Show that for all $x,y\in B$ we have $\mul[G](x,y)\in B$.
    \begin{subproof}
        Fix $x,y\in B$.
        Take $c\in\carrier[R]$ such that
            $x=\ringMul[R](a,c)$ by \auto.
        Take $d\in\carrier[R]$ such that
            $y=\ringMul[R](a,d)$ by \auto.
        $\ringAdd[R](c,d)\in\carrier[R]$ by \obv{ring_add_type}.
        $\ringAdd[R](x,y)
            =\ringMul[R](a,\ringAdd[R](c,d))$ by
            \auto{ring_distributive_left}.
        Follows by \auto.
    \end{subproof}

    Show that for every $x\in B$ we have $\groupInv[G](x)\in B$.
    \begin{subproof}
        Fix $x\in B$.
        Take $c\in\carrier[R]$ such that
            $x=\ringMul[R](a,c)$ by \auto.
        $\ringNeg[R](c)\in\carrier[R]$ by \obv{ring_neg_type}.
        $\ringNeg[R](x)=\ringMul[R](a,\ringNeg[R](c))$ by
            \auto{ring_neg_product_right}.
        Follows by \auto.
    \end{subproof}
    Hence $\subgroup{B}{G}$ by \auto{subgroup}.

    Show that for all $r,x$ such that $r\in\carrier[R]$ and $x\in B$
        we have $\ringMul[R](r,x)\in B$.
    \begin{subproof}
        Fix $r,x$.
        Assume $r\in\carrier[R]$ and $x\in B$.
        Take $c\in\carrier[R]$ such that
            $x=\ringMul[R](a,c)$ by \auto.
        $\ringMul[R](r,c)\in\carrier[R]$ by \obv{ring_mul_type}.
        $\ringMul[R](r,x)
            =\ringMul[R](a,\ringMul[R](r,c))$ by
            \auto{ring_mul_associative,
                commutative_ring_mul_commutative}.
        Follows by \auto.
    \end{subproof}
    Hence $B$ is a left ideal of $R$ by
        \auto{left_ideal_expanded_iff}.
    Hence $B$ is an ideal of $R$ by
        \auto{commutative_ideal_iff_left_ideal}.

    $\ringOne[R]\in\carrier[R]$ by \obv{ring_one_type}.
    $a=\ringMul[R](a,\ringOne[R])$ by \obv{ring_mul_one_right}.
    Hence $a\in B$ by \auto.
    Hence $\principalIdeal{R}{a}\subseteq B$ by
        \auto{principal_ideal_universal_property}.

    $\principalIdeal{R}{a}$ is an ideal of $R$ by
        \obv{principal_ideal_is_ideal}.
    Hence $\principalIdeal{R}{a}$ is a right ideal of $R$ by
        \auto{ideal,two_sided_ideal}.
    $a\in\principalIdeal{R}{a}$ by
        \obv{principal_ideal_contains_generator}.
    Show that $B\subseteq\principalIdeal{R}{a}$.
    \begin{subproof}
        Show that every element of $B$ is an element of
            $\principalIdeal{R}{a}$.
        \begin{subproof}
            Fix $x\in B$.
            Take $c\in\carrier[R]$ such that
                $x=\ringMul[R](a,c)$ by \auto.
            $\ringMul[R](a,c)\in\principalIdeal{R}{a}$ by
                \auto{right_ideal_expanded_iff}.
        \end{subproof}
        Follows by \obv{subseteq}.
    \end{subproof}
    Hence $\principalIdeal{R}{a}=B$ by
        \obv{subseteq_antisymmetric}.
    Follows by \auto.
\end{proof}

\begin{theorem}\label{principal_ideal_membership_iff_divisibility}
    Let $R$ be a commutative ring and $a,x\in\carrier[R]$.
    Then $x\in\principalIdeal{R}{a}$ iff $a$ divides $x$ in $R$.
\end{theorem}
\begin{proof}
    $R$ is a ring by \obv{commutative_ring_is_ring}.
    $\MultiplicativeMonoid{R}$ is a monoid by
        \obv{multiplicative_monoid_is_monoid}.
    Let $B=\{\ringMul[R](a,c)\mid c\in\carrier[R]\}$.
    $\principalIdeal{R}{a}=B$ by
        \obv{principal_ideal_eq_multiple_carrier}.

    Show that if $x\in\principalIdeal{R}{a}$, then
        $a$ divides $x$ in $R$.
    \begin{subproof}
        Assume $x\in\principalIdeal{R}{a}$.
        Hence $x\in B$.
        Take $c\in\carrier[R]$ such that
            $x=\ringMul[R](a,c)$ by \auto.
        $a,x,c\in\carrier[\MultiplicativeMonoid{R}]$ and
            $x=\mul[\MultiplicativeMonoid{R}](a,c)$ by
            \auto{multiplicative_monoid_fields}.
        Hence $a$ left-divides $x$ in $\MultiplicativeMonoid{R}$ by
            \auto{left_monoid_divisibility}.
        Hence $a$ left-divides $x$ in the ring $R$ by
            \auto{left_ring_divisibility}.
        Follows by \auto{ring_divisibility}.
    \end{subproof}

    Show that if $a$ divides $x$ in $R$, then
        $x\in\principalIdeal{R}{a}$.
    \begin{subproof}
        Assume $a$ divides $x$ in $R$.
        Hence $a$ left-divides $x$ in the ring $R$ by
            \obv{ring_divisibility}.
        Hence $a$ left-divides $x$ in $\MultiplicativeMonoid{R}$ by
            \obv{left_ring_divisibility}.
        Take $c\in\carrier[\MultiplicativeMonoid{R}]$ such that
            $x=\mul[\MultiplicativeMonoid{R}](a,c)$ by
            \auto{left_monoid_divisibility}.
        Hence $c\in\carrier[R]$ and $x=\ringMul[R](a,c)$ by
            \auto{multiplicative_monoid_fields}.
        Hence $x\in B$ by \auto.
        Follows by \auto.
    \end{subproof}
    Follows by \auto.
\end{proof}

\begin{theorem}\label{principal_ideal_power_containment}
    Let $R$ be a commutative ring and $a,d\in\carrier[R]$.
    Let $n$ be a natural number.
    Suppose $d\in\principalIdeal{R}{a}$.
    Then
        $\ringPower{R}{d}{n}
            \in\principalIdeal{R}{\ringPower{R}{a}{n}}$.
\end{theorem}
\begin{proof}
    $R$ is a ring by \obv{commutative_ring_is_ring}.
    Let $B=\{\ringMul[R](a,c)\mid c\in\carrier[R]\}$.
    $\principalIdeal{R}{a}=B$ by
        \obv{principal_ideal_eq_multiple_carrier}.
    Hence $d\in B$.
    Take $c\in\carrier[R]$ such that
        $d=\ringMul[R](a,c)$ by \auto.

    $\ringPower{R}{a}{n},\ringPower{R}{c}{n}\in\carrier[R]$ by
        \obv{ring_power_type}.
    $\ringPower{R}{d}{n}
        =\ringMul[R](\ringPower{R}{a}{n},
            \ringPower{R}{c}{n})$ by
        \auto{commutative_ring_power_of_product}.
    $\principalIdeal{R}{\ringPower{R}{a}{n}}
        =\{\ringMul[R](\ringPower{R}{a}{n},e)
            \mid e\in\carrier[R]\}$ by
        \obv{principal_ideal_eq_multiple_carrier}.
    Follows by \auto.
\end{proof}

\subsection{Annihilators of module subsets}

\begin{definition}\label{module_annihilator}
    Let $R$ be a ring.
    Let $M$ be a left module over $R$.
    Let $A$ be a subset of $\carrier[M]$.
    $\moduleAnnihilator{R}{M}{A}=
        \{a\in\carrier[R]\mid
            \text{for every $x\in A$ we have
                $\scalarMul[M](a,x)=\neutral[M]$}\}$.
\end{definition}

\begin{proposition}\label{module_annihilator_iff}
    Let $R$ be a ring.
    Let $M$ be a left module over $R$.
    Let $A$ be a subset of $\carrier[M]$.
    Then $a\in\moduleAnnihilator{R}{M}{A}$ iff
        $a\in\carrier[R]$ and
        for every $x\in A$ we have
            $\scalarMul[M](a,x)=\neutral[M]$.
\end{proposition}
\begin{proof}
    Follows by \obv{module_annihilator}.
\end{proof}

\begin{proposition}\label{module_annihilator_subseteq_carrier}
    Let $R$ be a ring.
    Let $M$ be a left module over $R$.
    Let $A$ be a subset of $\carrier[M]$.
    Then $\moduleAnnihilator{R}{M}{A}\subseteq\carrier[R]$.
\end{proposition}
\begin{proof}
    Follows by \obv{subseteq,module_annihilator_iff}.
\end{proof}

\begin{proposition}\label{module_annihilator_zero}
    Let $R$ be a ring.
    Let $M$ be a left module over $R$.
    Let $A$ be a subset of $\carrier[M]$.
    Then $\ringZero[R]\in\moduleAnnihilator{R}{M}{A}$.
\end{proposition}
\begin{proof}
    $\ringZero[R]\in\carrier[R]$ by \obv{ring_zero_type}.
    For every $x\in A$ we have $x\in\carrier[M]$ by \obv{subseteq}.
    Hence for every $x\in A$ we have
        $\scalarMul[M](\ringZero[R],x)=\neutral[M]$ by
        \obv{module_scalar_zero}.
    Follows by \obv{module_annihilator_iff}.
\end{proof}

\begin{proposition}\label{module_annihilator_addition}
    Let $R$ be a ring.
    Let $M$ be a left module over $R$.
    Let $A$ be a subset of $\carrier[M]$.
    Let $a,b\in\moduleAnnihilator{R}{M}{A}$.
    Then $\ringAdd[R](a,b)\in\moduleAnnihilator{R}{M}{A}$.
\end{proposition}
\begin{proof}
    $a,b\in\carrier[R]$ by \obv{module_annihilator_iff}.
    Hence $\ringAdd[R](a,b)\in\carrier[R]$ by \obv{ring_add_type}.
    Show that for every $x\in A$ we have
        $\scalarMul[M](\ringAdd[R](a,b),x)=\neutral[M]$.
    \begin{subproof}
        Fix $x\in A$.
        Hence $x\in\carrier[M]$ by \obv{subseteq}.
        $\scalarMul[M](a,x)=\neutral[M]$ and
            $\scalarMul[M](b,x)=\neutral[M]$ by
            \obv{module_annihilator_iff}.
        $\scalarMul[M](\ringAdd[R](a,b),x)
            =\mul[M](\scalarMul[M](a,x),\scalarMul[M](b,x))$ by
            \obv{module_scalar_add_distributive}.
        $\neutral[M]\in\carrier[M]$ by \obv{unitalmagma_type}.
        $\mul[M](\neutral[M],\neutral[M])=\neutral[M]$ by
            \obv{unitalmagma_left}.
    \end{subproof}
    Follows by \obv{module_annihilator_iff}.
\end{proof}

\begin{proposition}\label{module_annihilator_negation}
    Let $R$ be a ring.
    Let $M$ be a left module over $R$.
    Let $A$ be a subset of $\carrier[M]$.
    Let $a\in\moduleAnnihilator{R}{M}{A}$.
    Then $\ringNeg[R](a)\in\moduleAnnihilator{R}{M}{A}$.
\end{proposition}
\begin{proof}
    $a\in\carrier[R]$ by \obv{module_annihilator_iff}.
    Hence $\ringNeg[R](a)\in\carrier[R]$ by \obv{ring_neg_type}.
    Show that for every $x\in A$ we have
        $\scalarMul[M](\ringNeg[R](a),x)=\neutral[M]$.
    \begin{subproof}
        Fix $x\in A$.
        Hence $x\in\carrier[M]$ by \obv{subseteq}.
        $\scalarMul[M](a,x)=\neutral[M]$ by
            \obv{module_annihilator_iff}.
        $\scalarMul[M](\ringNeg[R](a),x)
            =\groupInv[M](\scalarMul[M](a,x))$ by
            \obv{module_scalar_negation}.
        Follows by \auto{module_is_group,group_inverse_neutral}.
    \end{subproof}
    Follows by \obv{module_annihilator_iff}.
\end{proof}

\begin{theorem}\label{module_annihilator_additive_subgroup}
    Let $R$ be a ring.
    Let $M$ be a left module over $R$.
    Let $A$ be a subset of $\carrier[M]$.
    Then
        $\subgroup{\moduleAnnihilator{R}{M}{A}}{\AdditiveGroup{R}}$.
\end{theorem}
\begin{proof}
    Let $I=\moduleAnnihilator{R}{M}{A}$ and $G=\AdditiveGroup{R}$.
    $G$ is a group by
        \auto{additive_group_is_exact_commutative_group,
            commutative_group_is_group}.
    $I\subseteq\carrier[G]$ by
        \auto{module_annihilator_subseteq_carrier,
            additive_group_fields}.
    $\neutral[G]\in I$ by
        \auto{additive_group_fields,module_annihilator_zero}.
    For all $a,b\in I$ we have $\mul[G](a,b)\in I$ by
        \auto{additive_group_fields,module_annihilator_addition}.
    For every $a\in I$ we have $\groupInv[G](a)\in I$ by
        \auto{additive_group_fields,module_annihilator_negation}.
    Follows by \auto{subgroup}.
\end{proof}

\begin{theorem}\label{module_annihilator_is_left_ideal}
    Let $R$ be a ring.
    Let $M$ be a left module over $R$.
    Let $A$ be a subset of $\carrier[M]$.
    Then $\moduleAnnihilator{R}{M}{A}$ is a left ideal of $R$.
\end{theorem}
\begin{proof}
    Let $I=\moduleAnnihilator{R}{M}{A}$.
    $\subgroup{I}{\AdditiveGroup{R}}$ by
        \auto{module_annihilator_additive_subgroup}.
    Show that for all $r,a$ such that $r\in\carrier[R]$ and $a\in I$
        we have $\ringMul[R](r,a)\in I$.
    \begin{subproof}
        Fix $r,a$.
        Assume $r\in\carrier[R]$ and $a\in I$.
        Hence $a\in\carrier[R]$ by \auto{module_annihilator_iff}.
        Hence $\ringMul[R](r,a)\in\carrier[R]$ by \obv{ring_mul_type}.
        Show that for every $x\in A$ we have
            $\scalarMul[M](\ringMul[R](r,a),x)=\neutral[M]$.
        \begin{subproof}
            Fix $x\in A$.
            Hence $x\in\carrier[M]$ by \obv{subseteq}.
            $\scalarMul[M](a,x)=\neutral[M]$ by
                \auto{module_annihilator_iff}.
            $\scalarMul[M](r,\scalarMul[M](a,x))
                =\scalarMul[M](\ringMul[R](r,a),x)$ by
                \obv{module_scalar_associative}.
            $\scalarMul[M](r,\neutral[M])=\neutral[M]$ by
                \obv{module_scalar_vector_zero}.
        \end{subproof}
        Follows by \obv{module_annihilator_iff}.
    \end{subproof}
    Follows by \auto{left_ideal_expanded_iff}.
\end{proof}

\begin{theorem}\label{module_annihilator_is_two_sided_ideal}
    Let $R$ be a ring.
    Let $M$ be a left module over $R$.
    Let $A$ be a subset of $\carrier[M]$.
    Suppose $A$ is a submodule of $M$ over $R$.
    Then $\moduleAnnihilator{R}{M}{A}$ is a two-sided ideal of $R$.
\end{theorem}
\begin{proof}
    $A\subseteq\carrier[M]$ by \obv{submodule_subseteq}.
    Let $I=\moduleAnnihilator{R}{M}{A}$.
    $I$ is a left ideal of $R$ by
        \auto{module_annihilator_is_left_ideal}.
    $\subgroup{I}{\AdditiveGroup{R}}$ by
        \auto{module_annihilator_additive_subgroup}.
    Show that for all $a,r$ such that $a\in I$ and $r\in\carrier[R]$
        we have $\ringMul[R](a,r)\in I$.
    \begin{subproof}
        Fix $a,r$.
        Assume $a\in I$ and $r\in\carrier[R]$.
        Hence $a\in\carrier[R]$ by \auto{module_annihilator_iff}.
        Hence $\ringMul[R](a,r)\in\carrier[R]$ by \obv{ring_mul_type}.
        Show that for every $x\in A$ we have
            $\scalarMul[M](\ringMul[R](a,r),x)=\neutral[M]$.
        \begin{subproof}
            Fix $x\in A$.
            $\scalarMul[M](r,x)\in A$ by \obv{submodule_scalar}.
            Hence
                $\scalarMul[M](a,\scalarMul[M](r,x))=\neutral[M]$ by
                \auto{module_annihilator_iff}.
            $x\in\carrier[M]$ by \obv{subseteq}.
            $\scalarMul[M](a,\scalarMul[M](r,x))
                =\scalarMul[M](\ringMul[R](a,r),x)$ by
                \obv{module_scalar_associative}.
        \end{subproof}
        Follows by \obv{module_annihilator_iff}.
    \end{subproof}
    Hence $I$ is a right ideal of $R$ by
        \auto{right_ideal_expanded_iff}.
    Follows by \auto{two_sided_ideal}.
\end{proof}

\subsection{Images and preimages}

\begin{theorem}\label{ring_hom_two_sided_ideal_preimage}
    Suppose $f$ is a ring homomorphism from $R$ to $S$.
    Suppose $J$ is a two-sided ideal of $S$.
    Then $\preimg{f}{J}$ is a two-sided ideal of $R$.
\end{theorem}
\begin{proof}
    Let $G=\AdditiveGroup{R}$ and $H=\AdditiveGroup{S}$.
    Let $P=\preimg{f}{J}$.
    $R$ is a ring and $S$ is a ring by
        \obv{ring_hom_source_ring,ring_hom_target_ring}.
    $\groupHom{f}{G}{H}$ by
        \auto{ring_hom_is_additive_group_hom}.
    $\subgroup{J}{H}$ by
        \auto{two_sided_ideal,left_ideal_additive_subgroup}.
    Hence $\subgroup{P}{G}$ by
        \auto{group_hom_subgroup_preimage}.
    Show that for all $a,x$ such that $a\in\carrier[R]$ and
        $x\in P$ we have $\ringMul[R](a,x)\in P$.
    \begin{subproof}
        Fix $a,x$.
        Assume $a\in\carrier[R]$ and $x\in P$.
        Hence $x\in\carrier[G]$ and $f(x)\in J$ by
            \auto{group_hom_preimage_iff}.
        Hence $x\in\carrier[R]$ by \auto{additive_group_fields}.
        $\ringMul[R](a,x)\in\carrier[R]$ by \obv{ring_mul_type}.
        Hence $\ringMul[R](a,x)\in\carrier[G]$ by
            \auto{additive_group_fields}.
        $f(a)\in\carrier[S]$ by \obv{ring_hom_apply_type}.
        $f(\ringMul[R](a,x))=\ringMul[S](f(a),f(x))$ by
            \obv{ring_hom_mul}.
        $\ringMul[S](f(a),f(x))\in J$ by
            \auto{two_sided_ideal,left_ideal_expanded_iff}.
        Hence $f(\ringMul[R](a,x))\in J$.
        Follows by \auto{group_hom_preimage_iff}.
    \end{subproof}
    Hence $P$ is a left ideal of $R$ by
        \auto{left_ideal_expanded_iff}.
    Show that for all $x,a$ such that $x\in P$ and
        $a\in\carrier[R]$ we have $\ringMul[R](x,a)\in P$.
    \begin{subproof}
        Fix $x,a$.
        Assume $x\in P$ and $a\in\carrier[R]$.
        Hence $x\in\carrier[G]$ and $f(x)\in J$ by
            \auto{group_hom_preimage_iff}.
        Hence $x\in\carrier[R]$ by \auto{additive_group_fields}.
        $\ringMul[R](x,a)\in\carrier[R]$ by \obv{ring_mul_type}.
        Hence $\ringMul[R](x,a)\in\carrier[G]$ by
            \auto{additive_group_fields}.
        $f(a)\in\carrier[S]$ by \obv{ring_hom_apply_type}.
        $f(\ringMul[R](x,a))=\ringMul[S](f(x),f(a))$ by
            \obv{ring_hom_mul}.
        $\ringMul[S](f(x),f(a))\in J$ by
            \auto{two_sided_ideal,right_ideal_expanded_iff}.
        Hence $f(\ringMul[R](x,a))\in J$.
        Follows by \auto{group_hom_preimage_iff}.
    \end{subproof}
    Hence $P$ is a right ideal of $R$ by
        \auto{right_ideal_expanded_iff}.
    Follows by \auto{two_sided_ideal}.
\end{proof}

\begin{theorem}\label{ring_surjection_two_sided_ideal_image}
    Suppose $f$ is a ring surjection from $R$ to $S$.
    Suppose $I$ is a two-sided ideal of $R$.
    Then $\img{f}{I}$ is a two-sided ideal of $S$.
\end{theorem}
\begin{proof}
    Let $G=\AdditiveGroup{R}$ and $H=\AdditiveGroup{S}$.
    Let $P=\img{f}{I}$.
    $f$ is a ring homomorphism from $R$ to $S$ and
        $f\in\surj{\carrier[R]}{\carrier[S]}$ by
        \obv{ring_surjection_is_hom,ring_surjection_is_surjection}.
    $R$ is a ring and $S$ is a ring by
        \obv{ring_hom_source_ring,ring_hom_target_ring}.
    $\groupHom{f}{G}{H}$ by
        \auto{ring_hom_is_additive_group_hom}.
    $\subgroup{I}{G}$ by
        \auto{two_sided_ideal,left_ideal_additive_subgroup}.
    Hence $\subgroup{P}{H}$ by
        \auto{group_hom_subgroup_image}.
    Show that for all $s,y$ such that $s\in\carrier[S]$ and
        $y\in P$ we have $\ringMul[S](s,y)\in P$.
    \begin{subproof}
        Fix $s,y$.
        Assume $s\in\carrier[S]$ and $y\in P$.
        Take $a\in\carrier[R]$ such that $f(a)=s$ by \obv{surj}.
        Take $x\in I$ such that $f(x)=y$ by
            \auto{group_hom_image_iff,subgroup_subseteq}.
        $\ringMul[R](a,x)\in I$ by
            \auto{two_sided_ideal,left_ideal_expanded_iff}.
        $x\in\carrier[R]$ by
            \auto{two_sided_ideal_subseteq_carrier,subseteq}.
        $f(\ringMul[R](a,x))=\ringMul[S](s,y)$ by
            \auto{ring_hom_mul}.
        Follows by
            \auto{group_hom_image_iff,subgroup_subseteq}.
    \end{subproof}
    Hence $P$ is a left ideal of $S$ by
        \auto{left_ideal_expanded_iff}.
    Show that for all $y,s$ such that $y\in P$ and
        $s\in\carrier[S]$ we have $\ringMul[S](y,s)\in P$.
    \begin{subproof}
        Fix $y,s$.
        Assume $y\in P$ and $s\in\carrier[S]$.
        Take $a\in\carrier[R]$ such that $f(a)=s$ by \obv{surj}.
        Take $x\in I$ such that $f(x)=y$ by
            \auto{group_hom_image_iff,subgroup_subseteq}.
        $\ringMul[R](x,a)\in I$ by
            \auto{two_sided_ideal,right_ideal_expanded_iff}.
        $x\in\carrier[R]$ by
            \auto{two_sided_ideal_subseteq_carrier,subseteq}.
        $f(\ringMul[R](x,a))=\ringMul[S](y,s)$ by
            \auto{ring_hom_mul}.
        Follows by
            \auto{group_hom_image_iff,subgroup_subseteq}.
    \end{subproof}
    Hence $P$ is a right ideal of $S$ by
        \auto{right_ideal_expanded_iff}.
    Follows by \auto{two_sided_ideal}.
\end{proof}

\subsection{Kernels and induced additive groups}

\begin{theorem}\label{ring_kernel_is_two_sided_ideal}
    Suppose $f$ is a ring homomorphism from $R$ to $S$.
    Then $\ringKernel{f}{R}$ is a two-sided ideal of $R$.
\end{theorem}
\begin{proof}
    Let $I=\ringKernel{f}{R}$.
    $R$ is a ring and $S$ is a ring by
        \obv{ring_hom_source_ring,ring_hom_target_ring}.
    $\subgroup{I}{\AdditiveGroup{R}}$ by
        \auto{ring_kernel_is_additive_subgroup}.
    Show that for all $a,x$ such that $a\in\carrier[R]$ and $x\in I$ we have
        $\ringMul[R](a,x)\in I$.
    \begin{subproof}
        Fix $a,x$.
        Assume $a\in\carrier[R]$ and $x\in I$.
        Hence $x\in\carrier[R]$ and $f(x)=\ringZero[S]$ by
            \auto{ring_kernel_iff}.
        $\ringMul[R](a,x)\in\carrier[R]$ by \obv{ring_mul_type}.
        $f(\ringMul[R](a,x))=\ringMul[S](f(a),f(x))$ by
            \obv{ring_hom_mul}.
        Hence $f(\ringMul[R](a,x))=\ringZero[S]$ by
            \auto{ring_hom_apply_type,ring_mul_zero_right}.
        Follows by \auto{ring_kernel_iff}.
    \end{subproof}
    Hence $I$ is a left ideal of $R$ by
        \auto{left_ideal_expanded_iff}.
    Show that for all $x,a$ such that $x\in I$ and $a\in\carrier[R]$ we have
        $\ringMul[R](x,a)\in I$.
    \begin{subproof}
        Fix $x,a$.
        Assume $x\in I$ and $a\in\carrier[R]$.
        Hence $x\in\carrier[R]$ and $f(x)=\ringZero[S]$ by
            \auto{ring_kernel_iff}.
        $\ringMul[R](x,a)\in\carrier[R]$ by \obv{ring_mul_type}.
        $f(\ringMul[R](x,a))=\ringMul[S](f(x),f(a))$ by
            \obv{ring_hom_mul}.
        Hence $f(\ringMul[R](x,a))=\ringZero[S]$ by
            \auto{ring_hom_apply_type,ring_mul_zero_left}.
        Follows by \auto{ring_kernel_iff}.
    \end{subproof}
    Hence $I$ is a right ideal of $R$ by
        \auto{right_ideal_expanded_iff}.
    Follows by \auto{two_sided_ideal}.
\end{proof}

\begin{theorem}\label{left_ideal_induced_additive_group}
    Suppose $I$ is a left ideal of $R$.
    Then
        $\ModuleAdditiveGroup{
            \InducedSubmodule{R}{\LeftRegularModule{R}}{I}}
        =\subgroupStructure{\AdditiveGroup{R}}{I}$.
\end{theorem}
\begin{proof}
    $I$ is a submodule of $\LeftRegularModule{R}$ over $R$ by
        \obv{left_ideal_is_submodule}.
    Hence
        $\ModuleAdditiveGroup{
            \InducedSubmodule{R}{\LeftRegularModule{R}}{I}}
        =\subgroupStructure{
            \ModuleAdditiveGroup{\LeftRegularModule{R}}}{I}$ by
        \obv{induced_submodule_additive_group}.
    $\ModuleAdditiveGroup{\LeftRegularModule{R}}
        =\AdditiveGroup{R}$ by
        \obv{left_regular_module_additive_group}.
\end{proof}
